% !TeX program = pdflatex
% Self-contained source: no external figures, bibliography, or input files.
% Compile twice with: pdflatex Arithmetic_Transcendence_Unified.tex
\documentclass[11pt,a4paper]{article}
\usepackage[T1]{fontenc}
\usepackage[utf8]{inputenc}
\usepackage[a4paper,margin=25mm,headheight=15pt,headsep=7mm,footskip=12mm]{geometry}
\usepackage{amsmath,amsthm,mathtools}
\usepackage{newtxtext,newtxmath}
\usepackage{microtype}
\usepackage{booktabs,longtable,array,tabularx}
\usepackage{enumitem}
\usepackage{fancyhdr}
\usepackage{xcolor}
\usepackage{listings}
\usepackage{hyperref}
\usepackage{bookmark}
\usepackage[nameinlink,noabbrev]{cleveref}
\hypersetup{hidelinks,pdftitle={Arithmetic Transcendence of Integrals, Quantiles, and Lemniscates},pdfauthor={KEIO},pdfsubject={A consolidated research manuscript with an audit of the transfer to Euler's constant},pdfkeywords={transcendence; E-functions; periods; beta quantiles; lemniscates; Euler constant}}
\pagestyle{fancy}
\fancyhf{}
\fancyhead[L]{\small\scshape Arithmetic transcendence}
\fancyhead[R]{\small\scshape Research manuscript}
\fancyfoot[C]{\small\thepage}
\renewcommand{\headrulewidth}{0.3pt}
\setlength{\parindent}{1.2em}
\setlength{\parskip}{3pt}
\linespread{1.06}
\setlength{\emergencystretch}{2.5em}
\setcounter{tocdepth}{2}
\setcounter{secnumdepth}{3}
\allowdisplaybreaks[1]
\numberwithin{equation}{section}
\newtheorem{theorem}{Theorem}[section]
\newtheorem{proposition}[theorem]{Proposition}
\newtheorem{lemma}[theorem]{Lemma}
\newtheorem{corollary}[theorem]{Corollary}
\newtheorem{conjecture}[theorem]{Conjecture}
\newtheorem{external}[theorem]{External theorem}
\theoremstyle{definition}
\newtheorem{definition}[theorem]{Definition}
\newtheorem{example}[theorem]{Example}
\theoremstyle{remark}
\newtheorem{remark}[theorem]{Remark}
\newcommand{\Q}{\mathbb Q}
\newcommand{\Qb}{\overline{\mathbb Q}}
\newcommand{\R}{\mathbb R}
\newcommand{\C}{\mathbb C}
\newcommand{\Z}{\mathbb Z}
\newcommand{\N}{\mathbb N}
\newcommand{\Ga}{\mathbb G_a}
\newcommand{\Gm}{\mathbb G_m}
\newcommand{\B}{\mathrm B}
\newcommand{\Ein}{\operatorname{Ein}}
\newcommand{\Jac}{\operatorname{Jac}}
\newcommand{\Lie}{\operatorname{Lie}}
\newcommand{\End}{\operatorname{End}}
\newcommand{\trdeg}{\operatorname{trdeg}}
\newcommand{\ord}{\operatorname{ord}}
\newcommand{\sgn}{\operatorname{sgn}}
\newcommand{\length}{\operatorname{length}}
\newcommand{\divisor}{\operatorname{div}}
\newcommand{\lcm}{\operatorname{lcm}}
\newcommand{\Span}{\operatorname{span}}
\newcommand{\im}{\operatorname{im}}
\newcommand{\dd}{\,\mathrm d}
\newcommand{\ii}{\mathrm i}
\newcommand{\e}{\mathrm e}
\newcommand{\abs}[1]{\left|#1\right|}
\newcolumntype{Y}{>{\raggedright\arraybackslash}X}
\newcolumntype{P}[1]{>{\raggedright\arraybackslash}p{#1}}
\lstset{basicstyle=\ttfamily\footnotesize,columns=fullflexible,keepspaces=true,breaklines=true,breakatwhitespace=false,showstringspaces=false,frame=single,framerule=0.2pt,xleftmargin=0.5em,xrightmargin=0.5em,aboveskip=10pt,belowskip=10pt,language=Python}
\crefname{theorem}{Theorem}{Theorems}
\Crefname{theorem}{Theorem}{Theorems}
\crefname{lemma}{Lemma}{Lemmas}
\Crefname{lemma}{Lemma}{Lemmas}
\crefname{proposition}{Proposition}{Propositions}
\Crefname{proposition}{Proposition}{Propositions}
\crefname{corollary}{Corollary}{Corollaries}
\Crefname{corollary}{Corollary}{Corollaries}
\crefname{external}{External theorem}{External theorems}
\Crefname{external}{External theorem}{External theorems}
\crefname{conjecture}{Conjecture}{Conjectures}
\Crefname{conjecture}{Conjecture}{Conjectures}
\crefname{section}{Section}{Sections}
\Crefname{section}{Section}{Sections}
\begin{document}
\begin{titlepage}
\centering
\vspace*{9mm}
{\small\scshape Consolidated English research manuscript\par}
\vspace{15mm}
{\LARGE\bfseries Arithmetic Transcendence of Integrals,\\[4pt] Quantiles, and Lemniscates\par}
\vspace{7mm}
{\large $E$-functions, cyclotomic periods, local logarithms,\\[3pt]
and the obstruction at Euler's constant\par}
\vspace{12mm}
{\large KEIO\par}
\vspace{3mm}
{\normalsize 11 October 2026\par}
\vspace{12mm}
\begin{minipage}{0.93\textwidth}
\textbf{Abstract.}
We assemble a sequence of investigations into the arithmetic nature of explicit
integrals, inverse distribution functions, and algebraic-curve lengths. An
antisymmetric iterated exponential integral is treated by differential-field
translations and $E$-function specialization. Reparametrized Gaussian integrals
give a finite classification of algebraic quantiles of truncated mixtures.
Logarithmic forms yield arithmetic classifications for Student distributions,
half-integral beta distributions, and shortest beta probability intervals.
A square-root cover turns polynomial lemniscate lengths into closed periods;
a cyclotomic homology argument treats rational arc divisions without requiring
simple Jacobians. Algebraic division ratios are reduced to torsion, and local
abelian logarithms combined with primitive divisors give rational-input and
bounded-degree results. Fermat-quotient torsion classifications, including a
generalized-Jacobian treatment of second-kind differentials, provide stability
under polynomial weights for a family of incomplete beta integrals. Finally,
we examine the attempted transfer to the Euler--Mascheroni constant. Explicit
coefficient locking, an essential-singularity obstruction, disjoint local
convergence domains, and a missing arithmetic zero-division principle explain
why the preceding results do not constitute a proof for that constant.

\medskip
\textbf{Status.} This is a unified proof manuscript, not a claim of established
publication priority or completed independent peer review. Classical inputs
are cited separately from the arguments assembled here. The paper does
\emph{not} prove the irrationality or transcendence of the Euler--Mascheroni
constant. Numerical certificates check computations, not the external
transcendence theorems or the logical correctness of every proof.

\medskip
\textbf{Keywords.} $E$-function; analytic subgroup theorem; generalized
Jacobian; incomplete beta function; quantile; lemniscate; primitive divisor;
Euler--Mascheroni constant.
\end{minipage}
\vfill
{\small Single-file \LaTeX{} edition. All formulas, tables, references, and the\par
representative verification listing are contained in this source.\par}
\end{titlepage}
\begingroup
\small
\setlength{\parskip}{0pt}
\linespread{1}\selectfont
\tableofcontents
\endgroup
\clearpage

\section{Introduction and logical organization}\label{sec:intro}
The common problem in this paper is to identify an arithmetic relation that
would have to exist if a specified real number were algebraic, and then to
exclude that relation by a theorem whose hypotheses can actually be checked.
The relevant relation is not the same in every example. For an $E$-function
value it is a functional algebraic relation; for a beta quantile it may be a
linear form in logarithms; for an abelian integral it can be a torsion condition;
for an arc length it is a closed-period representation. Confusing these
mechanisms is especially dangerous when attempting to pass to a limit.

The organization is thematic rather than chronological. Earlier special cases
are retained where they explain a hypothesis, but later extensions are stated
only once. In particular, the prime-leaf argument is subsumed by the
cyclotomic argument for all integral exponents, and the first-kind beta
argument is subsumed by its second-kind extension. An explicit concordance
with the source research notes is provided in \cref{app:concordance}.

\subsection{Principal families and their scope}
\begin{longtable}{@{}P{0.25\textwidth}P{0.47\textwidth}P{0.22\textwidth}@{}}
\toprule
Family & Conclusion developed in this manuscript & Principal input\\
\midrule
\endhead
Antisymmetric $\Ein$ integrals & Simultaneous algebraic independence, including
an additional iterated-integral coordinate & Differential algebra and
Siegel--Shidlovskii\\
Truncated Gaussian mixtures & At most $5m^2$ candidate locations for all
algebraic location--probability pairs & Beukers specialization\\
Student and half-integral beta laws & Parity classification and a three-pair
exception criterion & Baker and Gelfond--Schneider\\
Shortest beta intervals & Integer/noninteger arithmetic dichotomy for rational
parameters of integral sum & Logarithmic forms and a monotonicity lemma\\
Polynomial lemniscates & Transcendence of length; separation from logarithms
unless the curve is a circle & Closed periods and algebraic groups\\
Cyclotomic arc divisions & Rational divisions classified for every exponent;
algebraic ratios constrained by torsion & Analytic subgroup theorem\\
Local and bounded-degree inputs & Uniform rational-input results and an
explicit height-independent large-exponent bound & Local logarithms and
primitive divisors\\
Fermat beta integrals & Linear independence stable under arbitrary polynomial
weights, with an exact algebraic exception & Torsion packets and generalized
Jacobians\\
Euler's constant & Exact identities and obstructions, but no arithmetic
classification of $\gamma$ & Mixed arithmetic Gevrey framework\\
\bottomrule
\end{longtable}

\subsection{Conventions and distinctions}
The field $\Qb$ is the algebraic closure of $\Q$ in $\C$. An algebraic real
number belongs to $\Qb\cap\R$. The phrase ``the quantile is transcendental''
always refers to the exact inverse-function value, not to a terminating
printed decimal. All real powers of positive numbers use their positive
real branches; complex logarithms are fixed on the relevant path.

A period is an integral over an integral or rational homology class. An abelian
logarithm is a vector in a complex Lie algebra, not merely its class modulo a
period lattice. Whenever individual coordinates are compared, actual paths
and actual logarithm vectors are fixed. Torsion means finite order in the
specified algebraic group; torsion of a projected point is not silently
identified with torsion in the original Jacobian.

The symbols $\gamma$ and $\Gamma$ have distinct meanings:
\begin{equation}\label{eq:gamma-def}
 \gamma=\lim_{n\to\infty}(H_n-\log n)=-\Gamma'(1),\qquad
 H_n=\sum_{k=1}^n\frac1k,\quad H_0=0.
\end{equation}
The differential transcendence of the function $\Gamma$ does not settle the
arithmetic of $\Gamma'(1)$. Also, linear independence, algebraic independence,
and the transcendence of a single number are different conclusions. They
are distinguished throughout.

\subsection{Proven inputs, manuscript claims, and conditional conclusions}
Theorems attributed to named authors are external inputs. The remaining
numbered statements have proofs or explicit reductions below; their inclusion
records the claims of this consolidated research draft, not an assertion of
independent validation or novelty. No theorem below about quantiles or curves
uses an unproved claim about $\gamma$. The only implications involving
unproved arithmetic separation principles are marked as conditional in
\cref{sec:gamma}.

Tamiozzo's Remarks 5.7 and 8.16 concern, respectively, division points and
lengths of polynomial lemniscates \cite{Tamiozzo}. The proposed arguments in
\cref{sec:lemniscate,sec:division} address those formulations. They should not
be confused with Erd\H{o}s's extremal-length problem. This manuscript does not
establish a historical first solution, and the absence of an identical
statement in a literature search is not evidence of priority.

\section{External inputs and a reusable separation principle}\label{sec:inputs}
\subsection{Arithmetic differential functions}
An $E$-function is an entire series $f(z)=\sum a_nz^n/n!$ with algebraic
coefficients, exponentially bounded conjugates and common denominators, and
a linear differential equation over $\Qb(z)$. In the applications below,
these conditions are checked by coefficient formulas and a displayed finite
first-order system.

\begin{external}[Siegel--Shidlovskii and Beukers]\label[external]{thm:SS}
Let a vector of $E$-functions satisfy $Y'=M(z)Y$ over $\Qb(z)$. At a nonzero
algebraic point which is not a pole of $M$, the transcendence degree of the
values equals that of the functions over $\Qb(z)$. Moreover, a homogeneous
algebraic relation among the values lifts to a functional relation of the
same degree. In particular, functional linear independence implies linear
independence of the values at such a point.
\end{external}
The version and the linear consequence used here are Theorems 1.1 and 1.3
and Corollary 1.4 of Beukers \cite{Beukers}.

\begin{external}[Kolchin--Ostrowski, antiderivative form]\label[external]{thm:KO}
Let a differential field of characteristic zero have algebraically closed
constant field $C$. In an extension with no new constants, suppose
$f_1',\ldots,f_m'$ belong to the base field. If the $f_i$ are algebraically
dependent over that field, some nonzero $C$-linear combination of them
belongs to it.
\end{external}
A convenient reference is Srinivasan \cite{Srinivasan}. The derivative
hypothesis is imposed at a single level; it is not applied indiscriminately
to iterated antiderivatives at different levels.

\subsection{Logarithms and algebraic groups}
We use Hermite--Lindemann in the form
$0\ne a\in\Qb\Rightarrow e^a\notin\Qb$, and Gelfond--Schneider in the form
$a^b\notin\Qb$ for algebraic $a\ne0,1$ and algebraic irrational $b$, with
any fixed logarithm of $a$. Baker's theorem says that an algebraic-coefficient
linear form in fixed logarithms of nonzero algebraic numbers is either zero
or transcendental \cite{Baker}.

\begin{external}[Analytic subgroup theorem]\label[external]{thm:AST}
Let $G/\Qb$ be a connected commutative algebraic group, let
$u\in\Lie(G)_{\C}$ satisfy $\exp_G(u)\in G(\Qb)$, and let $V$ be a
$\Qb$-linear subspace with $u\in V_{\C}$. There is a connected algebraic
subgroup $H\subseteq G$ such that
\[
 u\in\Lie(H)_{\C}\subseteq V_{\C}.
\]
The statement includes logarithms whose exponential is the identity.
\end{external}
We use the strengthened formulation in Huber--W\"ustholz,
Theorem 6.2 \cite{HW}. A related consequence is that a nonzero closed period
of an algebraic rational differential on a smooth connected projective curve
is transcendental; the cycle must avoid its poles. See \cite[Theorem 13.9 and
Corollary 13.13]{HW}.

\begin{lemma}[Abelian versus affine logarithms]\label[lemma]{lem:affine-separation}
Let $A/\Qb$ be an abelian variety, let $u$ be a logarithm of an algebraic
point, and let $\lambda$ be an algebraic invariant differential. If
$\lambda(u)\ne0$, then
\begin{equation}\label{eq:affine-separation}
 \lambda(u)\notin\Qb+
 \Span_{\Qb}\{z\in\C:e^z\in\Qb^{\times}\}.
\end{equation}
\end{lemma}
\begin{proof}
Suppose $\lambda(u)=b_0+\sum b_jz_j$ with algebraic $b_j$ and algebraic
$e^{z_j}$. Apply \cref{thm:AST} to
$A\times\Ga\times\Gm^r$, the logarithm $(u,1,z_1,\ldots,z_r)$, and the
corresponding hyperplane. A connected algebraic subgroup of $A\times L$,
where $L$ is connected affine and commutative, is a product $B\times D$.
Indeed, its connected affine part maps trivially to $A$; after dividing by
that part, its abelian quotient maps trivially to the affine quotient of
$L$. These are the standard consequences of the Barsotti--Chevalley theorem
\cite{Milne}. Consequently $(u,0,\ldots,0)$ belongs to the same subgroup
Lie algebra. The hyperplane then gives $\lambda(u)=0$, a contradiction.
\end{proof}
In particular a nonzero algebraic linear form in coordinates of abelian
logarithms cannot be algebraic. Products of abelian varieties are allowed in
this assertion.

\subsection{Other inputs used later}
The geometric torsion packets used in \cref{sec:ratios,sec:fermat} are those
of Boxall--Grant \cite{BoxallGrant} and Arul \cite{Arul}; their exact special
cases are stated where applied. Manin--Mumford is used only for finiteness,
not for an effective enumeration \cite{Roessler}. The local logarithm is a
functorial homomorphism on a finite extension of $\Q_\ell$, kills torsion,
and agrees on a sufficiently small good-reduction residue disk with
termwise integration \cite{Zarhin}. Bang--Zsigmondy and Voutier provide
primitive divisors in \cref{sec:local}; their exceptional cases and numerical
thresholds are kept explicit.

\section{Antisymmetric iterated exponential integrals}\label{sec:ein}
Define the entire function
\begin{equation}
 \Ein(z)=\int_0^z\frac{1-e^{-t}}t\dd t
 =\sum_{n\ge1}\frac{(-1)^{n+1}z^n}{n\,n!}.
\end{equation}
For positive integers $i<j$, put
\begin{equation}\label{eq:Hij}
 F_i(z)=\Ein(iz),\qquad
 H_{ij}(z)=\int_0^z\bigl(F_j(t)F_i'(t)-F_i(t)F_j'(t)\bigr)\dd t.
\end{equation}
The integrands are entire, so these primitives are path independent.
The antisymmetry matters: the analogous sum integrates to $F_iF_j$ and
produces no additional coordinate.

\begin{theorem}\label{thm:ein}
For every $m\ge2$ and every $0\ne\alpha\in\Qb$, the numbers
\[
 e^{-\alpha},\quad F_i(\alpha)\ (1\le i\le m),\quad
 H_{ij}(\alpha)\ (1\le i<j\le m)
\]
are algebraically independent over $\Qb$. Their number is
$1+m+\binom m2$.
\end{theorem}
\begin{proof}
Set $u=e^{-z}$ and $K=\C(z,u)$, with derivation
$D=\partial_z-u\partial_u$. The exponential is transcendental over
$\C(z)$, as follows by comparing the largest exponential along the negative
real axis in a hypothetical polynomial relation. Thus $K$ is the rational
function field $\C(u)(z)$.

For $R\in K$, the derivative $DR$ cannot have exactly a simple pole at the
formal divisor $z=0$. If
$R=r_{-a}(u)z^{-a}+\cdots$, $a\ge1$, then its derivative has leading term
$-a r_{-a}(u)z^{-a-1}$. If $R$ is regular there, so is $DR$. This is a
valuation with $u$ treated as an independent coordinate, not the analytic
limit along $u=e^{-z}$.

Since $F_i'=(1-u^i)/z$, \cref{thm:KO} would turn algebraic dependence of
the $F_i$ over $K$ into
\[
 DR=\frac{\sum_i c_i(1-u^i)}z,
 \qquad (c_i)\ne0.
\]
The numerator is a nonzero polynomial. The simple pole contradicts the
preceding observation. Hence the $F_i$ are algebraically independent.

Let $L=K(F_1,\ldots,F_m)$. All $H_{ij}'$ lie in $L$. If the $H_{ij}$ were
algebraically dependent over $L$, another application of \cref{thm:KO}
would give
\[
 S=\sum_{i<j}c_{ij}H_{ij}\in L
\]
with not all $c_{ij}$ zero. For each $k$, let $\sigma_k$ translate
$F_k$ by $1$ and fix the other generators. These are commuting differential
automorphisms of $L$, and differentiation gives
\[
 \sigma_kS-S=\sum_{i<k}c_{ik}F_i-\sum_{j>k}c_{kj}F_j+d_k,
 \qquad d_k\in\C.
\]
For $i<j$, applying the two translations in opposite orders yields
\[
 (\sigma_i\sigma_j-\sigma_j\sigma_i)S=2c_{ij}.
\]
Commutativity forces every $c_{ij}$ to vanish. Therefore $u$, all $F_i$,
and all $H_{ij}$ are algebraically independent over $\C(z)$.

To specialize, write $H_{ij}(z)=\sum a_n^{ij}z^n/n!$. Direct convolution
shows
\begin{equation}\label{eq:ein-coeff}
 a_n^{ij}=\frac{(-1)^n}{n}\sum_{k=1}^{n-1}
 \binom nk\frac{j^ki^{n-k}-i^kj^{n-k}}k.
\end{equation}
If $L_n=\lcm(1,\ldots,n)$, then
\[
 |a_n^{ij}|\le2(i+j)^n,\qquad
 \operatorname{den}(a_n^{ij})\mid L_n^2.
\]
The standard estimate $\log L_n=O(n)$ gives the required exponential
common-denominator bound. A finite differential system is obtained by
including $u^r$ for $0\le r\le2m$, the $F_i$, the products $u^kF_i$,
and all $H_{ij}$; for example
\[
 (u^kF_i)'=-ku^kF_i+(u^k-u^{k+i})/z.
\]
Its only finite pole is $0$. All coordinates are $E$-functions and the
auxiliary coordinates add no transcendence degree. Apply
\cref{thm:SS} at $\alpha$.
\end{proof}

For one additional primitive, the implication ``algebraic primitive implies
it lies in the base field'' can also be seen by taking the normalized trace:
if $H'\in L$ and $H$ is algebraic of degree $d$, then
$R=d^{-1}\operatorname{Tr}(H)$ satisfies $R'=H'$. No new constants are
introduced in the meromorphic function fields used above.

\subsection{A distinguished constant and an infinite base field}
Set
\begin{equation}
 \mathcal A=H_{12}(1)
 =0.0668980772116667618622353744775514265199\ldots.
\end{equation}
In particular $e,\Ein(1),\Ein(2),\mathcal A$ are algebraically independent.
Since every polynomial involves only finitely many generators,
\cref{thm:ein} also gives
\begin{equation}
 \mathcal A\text{ is transcendental over }
 \Qb(e,\Ein(1),\Ein(2),\Ein(3),\ldots).
\end{equation}

For $F=F_1$, $G=F_2$, put $B=G-F$, $U=uG$, $V=u^2F$, and $H=H_{12}$.
The eight functions $(1,u,u^2,u^3,B,U,V,H)$ form the system
\begin{align}
 (u^r)'&=-ru^r\quad(1\le r\le3), & B'&=(u-u^2)/z,\\
 U'&=-U+(u-u^3)/z, & V'&=-2V+(u^2-u^3)/z,\\
 H'&=(B-U+V)/z. &&
\end{align}
They are linearly independent over $\Q(z)$ by the already established
algebraic independence. The classical quantitative linear-independence
estimate recalled in \cite{FRLiouville}, with dimension eight, gives
\begin{equation}
 2\le\mu(\mathcal A)\le8,
\end{equation}
where $\mu$ is the irrationality exponent. This is an upper bound, not an
assertion that the exponent equals eight; it also excludes Liouville behavior.

\subsection{An exact, but insufficient, connection to Euler's constant}
For $r,R>0$ define $H_r(R)$ as in \eqref{eq:Hij}, with scales $1,r$.
Differentiation in $r$, followed by integration by parts, gives
\begin{equation}
 r\partial_rH_r(R)=(1+e^{-rR})\Ein(R)
 +2\Ein(rR)-2\Ein((r+1)R).
\end{equation}
Using $\Ein(R)=\gamma+\log R+E_1(R)$ and $H_1=0$, integration in $r$
yields the locally uniform limit
\begin{equation}\label{eq:ein-gamma-limit}
 \lim_{R\to\infty}\{H_r(R)-\log r\log R\}
 =\gamma\log r-2\operatorname{Li}_2(-1/r)-\frac{\pi^2}{6}.
\end{equation}
The standard exponential-integral identity is in \cite{DLMF}. This limit
connects the iterated integral to familiar constants, but neither
transcendence at finite algebraic points nor algebraic independence is
automatically preserved by this limiting operation. Even transcendence of
the displayed limit would not by itself isolate $\gamma$ from the dilogarithm
and $\pi^2$ terms.

\section{Gaussian endpoint separation and finite quantile classification}\label{sec:gaussian}
For a positive integer $d$ put
\[
 J_d(x)=\int_0^x e^{-t^{2d}}\dd t.
\]
This is odd and strictly increasing on $\R$. To use $E$-functions one must
not simply call $J_d(z)$ an $E$-function. The appropriate reparametrization is
\begin{equation}
 E_{c,d}(z)=\int_0^1e^{-czt^{2d}}\dd t
 =\sum_{n\ge0}\frac{(-c)^n}{2dn+1}\frac{z^n}{n!},\qquad c>0.
\end{equation}
It satisfies
\begin{equation}\label{eq:gauss-ode}
 2dzE_{c,d}'+E_{c,d}=e^{-cz}.
\end{equation}
For algebraic $c$, conjugates of its normalized coefficients are exponentially
bounded. If $Mc$ is integral, a common denominator through order $n$ divides
$M^n\lcm(1,\ldots,2dn+1)$, which has exponential size.

\begin{lemma}\label[lemma]{lem:gauss-LI}
For distinct positive algebraic $r_1,\ldots,r_s$, the numbers
\[
 1,\ J_d(r_1),\ldots,J_d(r_s),\
 e^{-r_1^{2d}},\ldots,e^{-r_s^{2d}}
\]
are linearly independent over $\Qb$.
\end{lemma}
\begin{proof}
We first prove functional linear independence of
$1,E_{c_i,d},e^{-c_i z}$ for $0<c_1<\cdots<c_s$.
In a hypothetical rational-function relation, take the largest $c$ that
occurs and let $z=-x\to-\infty$. If its $E_{c,d}$ coefficient is zero, the
largest exponential cannot cancel. Otherwise the relation produces a rational
$R$ such that both $T=e^{cz}E_{c,d}-R$ and $T'$ decay exponentially, up to
polynomial factors. By \eqref{eq:gauss-ode}, the rational function
\[
 1-\{2dzR'+(1-2dcz)R\}
\]
is exponentially small and therefore identically zero.

But $2dzR'+(1-2dcz)R=1$ has no rational solution. At a finite nonzero pole,
the derivative raises the pole order. At a pole of order $a$ at zero the
leading coefficient is multiplied by $1-2da\ne0$. Thus $R$ would be a
polynomial, and then $-2dczR$ has an uncancellable leading degree.

The functions satisfy a finite $E$-system with only the finite singularity
$0$. Beukers's linear specialization theorem at $1$, and
$J_d(r)=rE_{r^{2d},d}(1)$, prove the assertion.
\end{proof}

\subsection{The endpoint-vector algorithm}
Let all data belong to a real number field $K$, and let
\begin{equation}\label{eq:mixture}
 f_d(x)=\sum_{j=1}^{m}w_j
 e^{-[\lambda_j(x-\mu_j)]^{2d}},\qquad a\le x\le b,
\end{equation}
where $a<b$ and $w_j,\lambda_j>0$. Merge identical $(\lambda_j,\mu_j)$.
A primitive is $S_d(x)=\sum_j(w_j/\lambda_j)J_d(\lambda_j(x-\mu_j))$.
For each positive radius $r$ define
\begin{equation}
 V_r(x)=\sum_{j:\,|\lambda_j(x-\mu_j)|=r}
 \frac{w_j}{\lambda_j}\sgn(x-\mu_j),
\end{equation}
omitting zero radii. By \cref{lem:gauss-LI}, for algebraic $x,q$,
\begin{equation}\label{eq:vector-criterion}
 \frac{S_d(x)-S_d(a)}{S_d(b)-S_d(a)}=q
 \quad\Longleftrightarrow\quad
 V(x)-(1-q)V(a)-qV(b)=0.
\end{equation}
The right side contains no transcendental values and is independent of $d$.

\begin{theorem}\label{thm:gauss-candidates}
All algebraic location--probability pairs in the interior of
\eqref{eq:mixture} lie in $K^2$ and can be found by inspecting at most $5m^2$
locations. Their set is independent of the positive integer $d$.
\end{theorem}
\begin{proof}
Let $\mathcal R$ contain all positive endpoint radii
$|\lambda_j(a-\mu_j)|$ and $|\lambda_j(b-\mu_j)|$; then
$|\mathcal R|\le2m$. The possible locations are the centers $\mu_j$, the
points $\mu_j\pm r/\lambda_j$ for $r\in\mathcal R$, and the solutions of
\[
 \lambda_j(x-\mu_j)=\varepsilon\lambda_k(x-\mu_k),
 \quad j<k,\quad\varepsilon=\pm1,
\]
with nonzero denominator. Their number is at most
$m+4m^2+2\binom m2=5m^2$. Outside this set every moving radius is nonzero,
distinct from the others, and distinct from the fixed endpoint radii. Its
coefficient in \eqref{eq:vector-criterion} cannot vanish.

Every candidate lies in $K$. Set $D=V(b)-V(a)\ne0$, the inequality following
from positivity of the total mass. A coordinate $D_r\ne0$ determines the
only possible value $q=(V_r(x)-V_r(a))/D_r\in K$. Test all remaining
coordinates and $0<q<1$. This also proves finiteness and independence of $d$.
The bound counts candidates, not bit operations.
\end{proof}

\subsection{Two explicit transcendental medians and genuine exceptions}
For
\[
 f(x)=e^{-(x+1/2)^2}+2e^{-(x-1/2)^2},\qquad -1\le x\le1,
\]
write $S(x)=J_1(x+1/2)+2J_1(x-1/2)$. The quantile equation becomes
\[
 S(x)=(-1+3q)J_1(1/2)+(-2+3q)J_1(3/2).
\]
For $-1<x<1$ no moving radius equals $3/2$. Thus an algebraic pair would
require $q=2/3$. The remaining candidate locations are $-1/2,0,1/2$; their
residuals are respectively
$-2J_1(1)-J_1(1/2)$, $-2J_1(1/2)$, and $J_1(1)-J_1(1/2)$, none zero.
All algebraic-probability quantiles of this distribution are therefore
transcendental, including its median
\[
 0.1111020192161737928704523032704820463312\ldots.
\]
Likewise the solution of $J_1(m_0)=J_1(1)/2$ is transcendental and equals
\[
 m_0=0.3926976257929472145472431237450194952618\ldots.
\]

On $[-3/2,3/2]$, however, the same asymmetric density has the exact pair
$(x,q)=(1/2,2/3)$, and the finite criterion shows it is the only interior
algebraic pair. A symmetric mixture on $[-1,1]$ has $(0,1/2)$; a single
Gaussian on $[0,1]$ has none. These examples prevent an unjustified statement
that all asymmetric distributions have exclusively transcendental quantiles.

\subsection{Why this does not settle a specified standard normal quantile}
For the untruncated standard normal distribution,
\[
 \Phi(x)=\frac12+\frac{J_1(x/\sqrt2)}{\sqrt\pi}.
\]
Two distinct positive algebraic quantiles at algebraic probabilities would
contradict \cref{lem:gauss-LI} after eliminating $\sqrt\pi$. Thus there is at
most one such positive algebraic quantile. This does not identify whether a
specified value such as $\Phi^{-1}(0.975)$ is that possible exception.
The normalization at infinity is a real obstruction to this particular proof.

\section{Logarithmic forms and statistical quantiles}\label{sec:classical-quantiles}
\subsection{A parity classification for Student distributions}
Let $Q_\nu(q)$ be the exact quantile of the standard Student distribution
with positive integral degrees of freedom $\nu$. The distribution definitions
and conventional tabulated critical values may be found in \cite{NIST}.

\begin{theorem}\label{thm:student}
For $q\in\Qb\cap(0,1)$ the following classification holds.
\begin{enumerate}[label=(\roman*),nosep]
\item If $\nu$ is even, $Q_\nu(q)$ is algebraic.
\item If $\nu\ge3$ is odd, $Q_\nu(q)$ is transcendental unless $q=1/2$,
in which case it is zero.
\item If $\nu=1$, the quantile is algebraic exactly when $q$ is rational.
\end{enumerate}
\end{theorem}
\begin{proof}
For $\nu=2m+1\ge3$, let $u=x/\sqrt\nu$. Reduction of the rational integral
gives
\begin{align}
 F_{2m+1}(x)&=\frac12+\frac{\arctan u+R_m(u)}\pi,\\
 R_m(u)&=\sum_{k=1}^m\frac{(2k-2)!!}{(2k-1)!!}
 \frac{u}{(1+u^2)^k}.
\end{align}
Every coefficient is positive, so $R_m(u)\ne0$ for real $u\ne0$. If $x,q$
were both algebraic, the quantile equation would give
\[
 \log\frac{1+\ii u}{1-\ii u}-(2q-1)\log(-1)=-2\ii R_m(u),
\]
with $\log(-1)=\ii\pi$ and the first logarithm equal to $2\ii\arctan u$.
The left side is a logarithmic form, whereas the right side is a nonzero
algebraic number. Baker's theorem excludes this equality.

For $\nu=2m$, set $y=x/\sqrt{2m+x^2}$. Then
\begin{equation}
 F_{2m}(x)=\frac12+c_m\sum_{k=0}^{m-1}
 (-1)^k\binom{m-1}{k}\frac{y^{2k+1}}{2k+1},\qquad
 c_m=\frac{(2m)!}{4^m m!(m-1)!}\in\Q.
\end{equation}
A specified algebraic value of this nonconstant polynomial forces $y$, and
hence $x=\sqrt{2m}\,y/\sqrt{1-y^2}$, to be algebraic.

Finally $Q_1(q)=\tan(\pi(q-1/2))$. Rational $q$ gives a root-of-unity
expression. For algebraic irrational $q$, algebraicity of the tangent would
make $(-1)^{2q-1}$ algebraic, contrary to Gelfond--Schneider.
\end{proof}

\begin{example}
At $q=39/40$, one has
\[
 Q_2(q)=19\sqrt{2/39},\qquad
 Q_3(q)=3.1824463052837095927232254257798683762689\ldots.
\]
The second number is transcendental. Here Hermite--Lindemann alone suffices:
if $u=Q_3(q)/\sqrt3$, $r=u/(1+u^2)$, and
$A=(1+\ii u)/(1-\ii u)$ were algebraic, then
\[
 \arctan u+r=19\pi/40\quad\Longrightarrow\quad
 A^{20}e^{40\ii r}=-1,
\]
an impossibility. The even-degree example $Q_4(q)$ instead satisfies
\[
 39X^6+468X^4-2928X^2-23104=0.
\]
These statements concern the exact values, not the rounded entries in a
table.
\end{example}

\subsection{Half-integral beta laws and three possible algebraic pairs}
For integers $n,m\ge0$, write
\[
 B_{n,m}(x)=\frac{1}{\B(n+1/2,m+1/2)}
 \int_0^x t^{n-1/2}(1-t)^{m-1/2}\dd t.
\]
When $n+m\ge1$, set
\[
 K_{n,m}=\frac{4^{n+m}n!m!(n+m)!}{(2n)!(2m)!}.
\]
There is a unique $P_{n,m}\in\Q[x]$ of degree at most $n+m-1$ satisfying
\begin{equation}\label{eq:half-beta-poly}
 2x(1-x)P'+(1-2x)P=2K_{n,m}x^n(1-x)^m-2.
\end{equation}
Indeed, the operator on the left is injective, raises the degree by one,
and has codimension-one image. Its image is the kernel of integration
against $[2\sqrt{x(1-x)}]^{-1}$. The right side has zero such integral by
normalization. Differentiating and checking the value at zero proves
\begin{equation}\label{eq:half-beta-primitive}
 B_{n,m}(x)=\frac{2\arcsin\sqrt x+
 \sqrt{x(1-x)}P_{n,m}(x)}\pi.
\end{equation}

\begin{theorem}\label{thm:three-pairs}
If $n+m\ge1$, and algebraic $x,q\in(0,1)$ satisfy $B_{n,m}(x)=q$, then
\begin{equation}\label{eq:three-pairs}
 (q,x)\in\left\{(1/3,1/4),(1/2,1/2),(2/3,3/4)\right\}.
\end{equation}
For $n,m\ge1$, at most one of these pairs occurs. A listed pair occurs
exactly when $P_{n,m}(x)=0$. If exactly one of $n,m$ is zero, none occurs.
\end{theorem}
\begin{proof}
For algebraic $x$, put $\theta=\arcsin\sqrt x$ and
$A=(1+\ii\sqrt{x/(1-x)})/(1-\ii\sqrt{x/(1-x)})$.
Equation \eqref{eq:half-beta-primitive} implies
\[
 q\log(-1)-\log A=\ii\sqrt{x(1-x)}P_{n,m}(x).
\]
Baker's theorem forces $P_{n,m}(x)=0$. Then $A=(-1)^q$ is algebraic, so
Gelfond--Schneider forces $q\in\Q$.

For $n,m\ge1$, let
$H(x)=B_{n,m}(x)-2\arcsin\sqrt x/\pi$. Its derivative is
\[
 H'(x)=\frac{K_{n,m}x^n(1-x)^m-1}{\pi\sqrt{x(1-x)}}.
\]
The numerator changes sign in the order $-,+,-$. Its positive maximum
exceeds zero because the corresponding density ratio has weighted mean one.
Since $H(0)=H(1)=0$, the function has exactly one zero in $(0,1)$.
Consequently $P_{n,m}$ has exactly one root there.

As $q$ is rational,
$x=(2-\zeta-\zeta^{-1})/4$ for the root of unity $\zeta=e^{\ii\pi q}$.
All algebraic conjugates of $x$ lie strictly in $(0,1)$, and all are roots
of the rational polynomial $P_{n,m}$. Uniqueness of its real interior root
therefore forces $x\in\Q$. Now
$2\cos(\pi q)=2-4x$ is a rational algebraic integer in $(-2,2)$, hence is
$-1,0$, or $1$. This gives \eqref{eq:three-pairs}.

If $n=0<m$, the density ratio $K_{0,m}(1-x)^m$ is strictly decreasing, and
$H$ is positive in the interior; for $m=0<n$ it is negative. Thus no root
occurs. Conversely, every pair in \eqref{eq:three-pairs} has
$2\arcsin\sqrt x=\pi q$, so $P(x)=0$ is sufficient.
\end{proof}
For $n=m=0$ the law is the arcsine distribution; its algebraic-probability
quantile is algebraic exactly when the probability is rational. This
excluded case cannot be folded into \cref{thm:three-pairs}.

A practical exact recurrence follows by writing the right side of
\eqref{eq:half-beta-poly} as $\sum h_kx^k$ and $P=\sum a_kx^k$:
\[
 a_{-1}=0,\qquad a_k=\frac{h_k+2ka_{k-1}}{2k+1},
 \qquad -2(n+m)a_{n+m-1}=h_{n+m}.
\]
This tests the three candidate locations for each specified parameter pair.
It does not classify all integer pairs $(n,m)$ for which $P_{n,m}(1/4)=0$.

\begin{example}
For $\operatorname{Beta}(3/2,5/2)$,
\[
 P_{1,2}(x)=-\frac23(4x-1)(2x-3).
\]
The only interior algebraic pair is $(q,x)=(1/3,1/4)$. The median and
both $95\%$ equal-tailed endpoints are transcendental:
\begin{align*}
 Q(1/2)&=0.3524522798771308137655997435743940518178\ldots,\\
 Q(1/40)&=0.0387476177851651607807019421865466897427\ldots,\\
 Q(39/40)&=0.8232639028687426825628811882915255744853\ldots.
\end{align*}
\end{example}

\subsection{Binomial intervals and the \texorpdfstring{$F$}{F} distribution}
With $s$ successes in $N\ge1$ Bernoulli observations, the Jeffreys posterior
is $\operatorname{Beta}(s+1/2,N-s+1/2)$. Thus its unmodified $95\%$
equal-tailed posterior interval has two transcendental endpoints, since
$1/40,39/40$ are absent from \eqref{eq:three-pairs}. Replacing an endpoint
by $0$ or $1$ in a boundary-adjusted implementation is a different definition
and of course gives an algebraic replacement. In contrast, the
Clopper--Pearson endpoints use integer beta parameters and rational
probabilities; their distribution functions are rational polynomials, so
all exact endpoints are algebraic. The statistical definitions are described
in \cite{Ratesci}. This arithmetic distinction asserts no coverage or
performance superiority.

For $Y$ with an $F(d_1,d_2)$ law,
$X=d_1Y/(d_1Y+d_2)$ has beta parameters $(d_1/2,d_2/2)$, and the inverse
transformation is rational. If one degree of freedom is even, all
algebraic-probability quantiles are algebraic, by a polynomial integral after
a square-root substitution if necessary. If both are odd and not both one,
there is at most one interior algebraic pair, at a probability in
$\{1/3,1/2,2/3\}$. In particular every such $95\%$ quantile is
transcendental. For $F(3,5)$ the sole exception is $Q(1/3)=5/9$.

\section{Shortest beta probability intervals}\label{sec:hpd}
The previous results prescribe a single probability $F(x)$. A shortest
interval instead imposes two simultaneous equations,
$F(b)-F(a)=c$ and $f(a)=f(b)$; neither $F(a)$ nor $F(b)$ need be algebraic.
This requires a new noncancellation argument.

\begin{theorem}\label{thm:hpd}
Let $\alpha,\beta\in\Q$ satisfy $\alpha,\beta>1$ and
$\alpha+\beta\in\Z$. For algebraic $c\in(0,1)$, let $[a,b]$ be the
unique shortest interval carrying beta probability $c$, measured in the
original $x$ coordinate. If the parameters are integers, $a,b,b-a$ are
algebraic. Otherwise each is transcendental, and
\[
 \trdeg_{\Qb}\Qb(a,b)=1.
\]
In the nonintegral case the midpoint is transcendental unless
$\alpha=\beta$, when it is $1/2$.
\end{theorem}

\subsection{Existence, uniqueness, and the density constraint}
For $\alpha,\beta>1$, the beta density is zero at both ends, strictly
increasing up to its mode, and then strictly decreasing. Its positive
superlevel sets are intervals $I_\lambda=[a_\lambda,b_\lambda]$ whose
mass runs continuously and strictly from zero to one. Thus the interval
of any specified mass exists uniquely and has equal endpoint densities.
If $J$ were shorter than $I_\lambda$, then
\[
 \int_{I_\lambda}f-\int_Jf
 \ge\lambda\bigl(|I_\lambda\setminus J|-|J\setminus I_\lambda|\bigr)
 =\lambda(|I_\lambda|-|J|)>0.
\]
This proves optimality, and strict unimodality gives uniqueness. It is the
usual highest-density interval, not an equal-tailed interval; the endpoint
condition also appears in \cite{Lourie}.

\subsection{A monotone correction divided by the density}
For nonintegral parameters write
\[
 \alpha=n+r,\quad\beta=m+1-r,\quad n,m\ge1,\quad0<r<1,
 \qquad r\in\Q.
\]
Set $M=n+m$ and
\[
 w_0=x^{r-1}(1-x)^{-r},\quad
 B_0=\B(r,1-r)=\frac\pi{\sin\pi r},\quad
 K=\frac{M!}{(r)_n(1-r)_m}.
\]
There is a unique rational polynomial $P$ of degree at most $M-1$ with
\begin{equation}\label{eq:hpd-P}
 x(1-x)P'+(r-x)P=Kx^n(1-x)^m-1.
\end{equation}
The same codimension-one weighted-integral argument as before proves this
assertion. If $D=x^r(1-x)^{1-r}P$, then
\begin{equation}\label{eq:hpd-decomp}
 B_0F(x)=\int_0^xw_0(t)\dd t+D(x).
\end{equation}

\begin{lemma}\label[lemma]{lem:hpd-monotone}
The function $R(x)=P(x)/[x^{n-1}(1-x)^{m-1}]$ satisfies $R'(x)>0$ on
$(0,1)$. Consequently $f(a)=f(b)>0$ and $a<b$ imply $D(b)-D(a)>0$.
\end{lemma}
\begin{proof}
Write $P=\sum_{k=0}^{M-1}p_kx^k(1-x)^{M-1-k}$. With
$p_{-1}=p_M=0$, coefficient comparison in \eqref{eq:hpd-P} yields
\[
 (j+r)p_j-(M-j+1-r)p_{j-1}
 =K\mathbf1_{j=n}-\binom Mj.
\]
The forward recurrence from $p_0=-1/r$ gives $p_k<0$ for $k<n$; the
backward recurrence from $p_{M-1}=1/(1-r)$ gives $p_k>0$ for $k\ge n$.
With $j=k-n+1$, a term of $R$ is $p_kx^j(1-x)^{1-j}$, whose derivative is
\[
 p_kx^{j-1}(1-x)^{-j}(j-x)>0.
\]
For $k<n$, both $p_k$ and $j-x$ are negative; otherwise both are positive.
Summing proves the assertion. Finally
\[
 D(x)=\frac{B_0}{K}f(x)R(x),
\]
so equal positive densities imply strict increase of $D$ between the two
endpoints.
\end{proof}

\subsection{Transcendence of endpoints, width, and midpoint}
Suppose $a,b$ were algebraic. By \cref{lem:hpd-monotone}, $D(b)-D(a)$ is a
nonzero algebraic number. But \eqref{eq:hpd-decomp} gives
\[
 D(b)-D(a)=cB_0-\int_a^bw_0(x)\dd x.
\]
Write $r=p/q$ in lowest terms and substitute
$y=(x/(1-x))^{1/q}$. Then
\[
 w_0(x)\dd x=\frac{qy^{p-1}}{1+y^q}\dd y.
\]
This is a proper rational function with simple algebraic poles. Its integral
between algebraic endpoints is an algebraic-coefficient logarithmic form.
The same is true of $cB_0$, since $\sin\pi r$ is nonzero algebraic.
Baker's theorem contradicts the nonzero algebraic left side.

Clearing denominators in the density equation gives
\begin{equation}\label{eq:hpd-curve}
 h(a)=h(b),\qquad h(X)=X^A(1-X)^B\in\Q[X],\quad A,B\in\Z_{>0}.
\end{equation}
If one endpoint were algebraic, the other would also be algebraic. Thus each
is transcendental, while their polynomial relation makes their joint
transcendence degree exactly one.

If $w=b-a>0$ were algebraic, $a$ would solve the nonzero polynomial
$h(X+w)-h(X)$ over $\Qb$. Its leading term is $d\kappa wX^{d-1}$ when
$\kappa X^d$ is the leading term of $h$. This is impossible. If the midpoint
$z=(a+b)/2$ were algebraic, the analogous polynomial
$h(2z-X)-h(X)$ would have to vanish identically. Reflection must interchange
the only roots $0,1$, preserving multiplicities, so $z=1/2$ and $A=B$.
This is exactly the symmetric case.

For integer parameters let $A=\alpha-1$, $B=\beta-1$. The equal-density
curve has the rational parametrization
\[
 a(t)=\frac{t^A-1}{t^{A+B}-1},\qquad b(t)=t^Ba(t),\qquad t>1.
\]
The mass $F(b(t))-F(a(t))$ is a nonconstant rational function, tending to
zero and one at the two ends of this parameter range. At an algebraic mass,
$t$ and hence the endpoints and width are algebraic. This completes the
proof of \cref{thm:hpd}.

\subsection{Examples and posterior interpretation}
For $\operatorname{Beta}(3/2,5/2)$ at mass $19/20$, the defining equations are
\[
 a(1-a)^3=b(1-b)^3,\qquad
 \frac{16}\pi\int_a^b\sqrt{x}(1-x)^{3/2}\dd x=\frac{19}{20}.
\]
The endpoints and width are
\begin{align*}
 a&=0.0095568708730449428602194557889816354395\ldots,\\
 b&=0.7707572692573354283348360180149011070037\ldots,\\
 b-a&=0.7612003983842904854746165622259194715642\ldots.
\end{align*}
All three, and their nonsymmetric midpoint, are transcendental. In the
parametrization $S=1+t+t^2+t^3$, $a=1/S$, $b=t^3/S$, one can independently
check
\[
 D(b)-D(a)=\frac{2(t-1)\sqrt{t(t^2+t+1)}}{3S^3}
 (3t^6+14t^5+25t^4+36t^3+25t^2+14t+3)>0.
\]
For the integer-parameter comparison $\operatorname{Beta}(2,3)$, the
$95\%$ value of $t$ instead solves
\[
 t^8+4t^7+10t^6-64t^5-361t^4-544t^3-390t^2-156t-39=0.
\]

If both successes and failures are at least one, a Jeffreys binomial
posterior has nonintegral half-integer parameters of integral sum. Its
shortest interval at every algebraic mass has transcendental endpoints and
width. A uniform-prior posterior has integer parameters, and its corresponding
quantities are algebraic. The conclusion is coordinate dependent: a shortest
interval after a nonlinear reparametrization is a different optimization
problem. Boundary data, U-shaped laws, and arbitrary rational parameter sums
are not covered by this theorem.

\section{Polynomial lemniscates: lengths as closed periods}\label{sec:lemniscate}
For nonconstant $P\in\Qb[z]$ and algebraic $\rho>0$, let
\[
 \Lambda_{P,\rho}=\{z\in\C:|P(z)|=\rho\},\qquad
 L(P,\rho)=\length(\Lambda_{P,\rho}).
\]
Length is that of the geometric set; repeated polynomial factors do not
multiply it. The argument below includes critical levels, crossing branches,
and multiple connected components.

\begin{theorem}\label{thm:lem-length}
The number $L(P,\rho)$ is transcendental. Moreover,
\begin{equation}\label{eq:circle-dichotomy}
 \frac{L(P,\rho)}{2\pi}\in\Qb
 \quad\Longleftrightarrow\quad
 P\text{ has exactly one distinct zero}.
\end{equation}
If $P$ has at least two distinct zeros, then $1,\pi,L(P,\rho)$ are linearly
independent over $\Qb$, and more generally $L(P,\rho)$ is not an algebraic
number plus an algebraic linear combination of logarithms of algebraic
numbers. In particular $e^{aL(P,\rho)}$ is transcendental for every
$0\ne a\in\Qb$.
\end{theorem}
This formulation extends the length question posed in
\cite[Remark 8.16]{Tamiozzo}. It makes no assertion about which polynomial
of fixed degree maximizes the length.

\subsection{Removing powers and constructing one connected curve}
Write $P/\rho=Q^d$ so that the multiplicities $m_j$ of the distinct zeros
of $Q$ have greatest common divisor one. A $d$th root of the leading
coefficient is algebraic, and the level set is unchanged. Put $n=\deg Q$.
Consider the complex curve
\begin{equation}\label{eq:complex-lem}
 Q(z)\overline Q(w)=1,
\end{equation}
where conjugation in $\overline Q$ acts on coefficients.

\begin{lemma}\label[lemma]{lem:lem-irreducible}
The curve \eqref{eq:complex-lem} is geometrically irreducible.
\end{lemma}
\begin{proof}
At a zero $a_j$ of $Q$ of multiplicity $m_j$, every root $\xi$ of
$\overline Q(w)-1/Q(z)$ has valuation $-m_j/n$. Nonnegative valuation
would be incompatible with the pole of $1/Q$; at negative valuation, the
leading term of $\overline Q$ is uniquely minimal.
If a monic factor of degree $e$ existed over $\Qb(z)$, its nonzero constant
coefficient, a product of $e$ such roots, would have integral valuation
$-em_j/n$. Thus $n\mid em_j$ for every $j$, and the greatest-common-divisor
condition forces $n\mid e$. There is no proper factor. Gauss's lemma gives
geometric irreducibility over the algebraically closed coefficient field.
\end{proof}
Related irreducibility statements occur in Orevkov \cite{Orevkov}; the
valuation argument is included to specify the connected model used here.
Let $X$ be its smooth projective normalization, with real structure
$(z,w)\mapsto(\overline w,\overline z)$. Set
$x=(z+w)/2$, $y=(z-w)/(2\ii)$. The real affine part is the original level
set. Its highest homogeneous real term is
$|a_n|^2(x^2+y^2)^n$, so it has no real point at infinity.

\subsection{Critical points are crossings, not cusps}
At a point $z_0$ of the level set, $Q(z_0)\ne0$, so locally
\[
 g(z)=\log\frac{Q(z)}{Q(z_0)}
\]
is holomorphic. If its zero has order $m$, choose a local holomorphic
coordinate $u$ with nonzero derivative so that $g=u^m$. The real level set
is then $\operatorname{Re}(u^m)=0$, the union of $m$ distinct straight lines.
Each line is a smooth immersed branch in the $z$ plane after reversing the
coordinate change. Normalization separates these branches. Hence on every
real branch of $X$,
\begin{equation}\label{eq:immersion}
 \left(\frac{dx}{dt}\right)^2+\left(\frac{dy}{dt}\right)^2>0.
\end{equation}
The compact real locus of $X$ is a finite disjoint union of circles. It is
nonempty, since a noncritical value of $Q$ on the unit circle has a regular
inverse branch. The finite collection of crossing points has length zero,
so summing the lengths of the normalized circles gives precisely the length
of the geometric level set.

\subsection{The positive square root closes on an algebraic cover}
On $X$ set $q=dz\,dw=(dx)^2+(dy)^2$. Choose an algebraic nonzero rational
one-form $\tau$ and write $q=f\tau^2$. If $f$ is not a square, take the
connected degree-two cover $\pi:Y\to X$ with function field obtained by
adjoining $v^2=f$, and put $\omega=v\pi^*\tau$. If $f$ is a square, use
$Y=X$ and a global square root. In either case
\begin{equation}\label{eq:sqrt-form}
 \omega^2=\pi^*q.
\end{equation}

The square root alone does not prove that the length is a closed period.
On an oriented real circle $C\subset X(\R)$, choose at each point the sheet
on which $\omega$ evaluates positively on the tangent direction.
By \eqref{eq:immersion} the cover is unramified there, and this choice is
unique and continuous. The point and oriented tangent agree after a complete
turn, so the chosen sheet also agrees. Thus the lifted path $\widetilde C$
is closed and
\[
 \int_{\widetilde C}\omega=\length(C\text{ in the plane})>0.
\]
If the square root is global, reverse the orientation of each circle when
necessary. All circles can then be combined on the same connected $Y$ into
an integral cycle $\gamma$ satisfying
\begin{equation}\label{eq:lem-period}
 \partial\gamma=0,\qquad \int_\gamma\omega=L(P,\rho)>0.
\end{equation}
The closed-period theorem quoted in \cref{sec:inputs} proves transcendence.
This is one nonzero period, not an unsupported sum-of-transcendentals
argument. No limiting passage from regular to singular levels is used.

\subsection{Global holomorphy and the exact circle exception}
Suppose $Q$ has at least two distinct zeros. The differential $q=dz\,dw$
has no pole where $z,w$ are finite. At an infinite point, exactly one of
$z,w$ tends to infinity and the other approaches a root of the other
polynomial. Suppose $z=\infty$, $w=a$, where $a$ has multiplicity $m<n$.
With $d=\gcd(n,m)$, normalization gives a local parameter $t$ with
\[
 z\sim c\,t^{-m/d},\qquad w-a\sim c'\,t^{n/d}.
\]
Consequently
\begin{equation}\label{eq:lem-order}
 \ord_t(q)=\frac{n-m}{d}-2\ge-1.
\end{equation}
The same holds with $z,w$ interchanged. If locally
$q=t^kh(t)(dt)^2$, $h(0)\ne0$, then for even $k$ the square root has order
$k/2\ge0$. For odd $k$, the substitution $t=u^2$ makes its order
$k+1\ge0$. Thus $\omega$ is globally holomorphic on $Y$.

The period \eqref{eq:lem-period} is now a coordinate of an abelian logarithm
of the identity in $\Jac(Y)$. Apply \cref{lem:affine-separation} to obtain
the stronger logarithmic separation in \cref{thm:lem-length}. If
$P(z)=c(z-a)^n$, the level set is a circle of algebraic radius
$(\rho/|c|)^{1/n}$ and its length is $2\pi$ times that radius. In this
case $m=n$ in the order calculation, producing a double pole of $q$ and a
simple pole of its square root. This proves the precise dichotomy and all
stated consequences.

\subsection{A necessary counterexample outside the hypothesis}
The algebraic curve parametrized by
$z(t)=(1+e^{\ii t})^2$, $-\pi\le t\le\pi$, has speed
$4\cos(t/2)$ and total length $16$. It has a cusp where the normalization
map is not an immersion. A positive square-root branch cannot be continued
through that point as the same closed holomorphic lift. This prevents an
incorrect extension of \cref{thm:lem-length} to every closed algebraic curve.

\subsection{A convergent certificate for a concrete normalized length}
Let $L_3=L(z^3-z,1)$. Define $y(t)$ by
$y^3-ty=1$, $y(0)=1$, and
\[
 G(t)=\sqrt{\frac3{3y(t)^2-t}}=\sum_{n\ge0}g_nt^n,\qquad G(0)=1.
\]
All $g_n$ are rational. The exterior parametrization
$\phi(u)=u y(u^{-2})$ satisfies $\phi^3-\phi=u^3$ and parametrizes the
level curve once on $|u|=1$. Its derivative is $G(u^{-2})^2$.
Fourier orthogonality therefore gives
\begin{equation}\label{eq:L3-series}
 \frac{L_3}{2\pi}=\sum_{n\ge0}g_n^2.
\end{equation}
The coefficients can be computed without numerical root finding:
\[
 y_0=1,\quad
 y_n=\frac{(-1)^{n-1}}{3n}\binom{(2n-4)/3}{n-1},\quad
 d_n=(1-2n)y_n,
\]
\[
 g_0=1,\qquad
 g_n=\frac12\left(d_n-\sum_{k=1}^{n-1}g_kg_{n-k}\right).
\]
On $|t|\le3/2$, all roots of $y^3-ty-1$ have modulus less than $3/2$.
Since
\[
 \prod_{j=1}^3(3y_j^2-t)=27-4t^3,
\]
we obtain $|3y_j^2-t|\ge24/121$ and $|G(t)|^2\le121/8$.
Cauchy's estimate gives the explicit tail bound
\begin{equation}\label{eq:L3-tail}
 \sum_{n>M}g_n^2\le\frac{1089}{40}\left(\frac49\right)^{M+1}.
\end{equation}
The values are
\begin{align*}
 L_3&=6.4647063278134772656231267328429160327830\ldots,\\
 L_3/\pi&=2.0577799354179392949673832718280913191325\ldots.
\end{align*}
Both are transcendental, and $1,\pi,L_3$ are linearly independent.
Algebraic independence of $\pi$ and $L_3$ is not asserted.

\section{Cyclotomic homology and all rational arc divisions}\label{sec:division}
For every positive integer $N$, put
\begin{equation}\label{eq:IN}
 I_N(s)=\int_0^s\frac{dt}{\sqrt{1-t^N}},\qquad
 R_N(s)=\frac{I_N(s)}{I_N(1)},\qquad0\le s\le1.
\end{equation}
Let $s_{N,q}$ be the unique solution of $R_N(s)=q$ in $(0,1)$.

\begin{theorem}\label{thm:all-division}
For every rational $q\in(0,1)$,
\[
 s_{N,q}\in\Qb\quad\Longleftrightarrow\quad N\in\{1,2,3,4,6\}.
\]
For every other positive integral exponent, each of these division values is
transcendental.
\end{theorem}
The distinction between the exponent $N$ and a lemniscate's number of leaves
is important: a curve with $n$ leaves uses $N=2n$.

\subsection{The primitive cyclotomic factor}
For $N\ge3$ consider $\mathcal H_N:y^2=1-x^N$ and its Jacobian $J_N$.
Its genus is $g=\lfloor(N-1)/2\rfloor$, with holomorphic basis
\[
 \omega_k=x^{k-1}dx/y,\qquad1\le k\le g.
\]
The rotation $T:(x,y)\mapsto(\zeta_Nx,y)$ has eigenvalue $\zeta_N^k$
on $\omega_k$. Its characteristic polynomial on rational homology is
\begin{equation}\label{eq:chiN}
 \chi_N(X)=\begin{cases}
 (X^N-1)/(X-1),&N\text{ odd},\\
 (X^N-1)/(X^2-1),&N\text{ even}.
 \end{cases}
\end{equation}
All eigenvalues occur once. A rational polynomial idempotent $e_N$ selects
the $\Phi_N$ factor:
\[
 e_N\equiv1\pmod{\Phi_N},\qquad
 e_N\equiv0\pmod{\chi_N/\Phi_N}.
\]
After clearing denominators by $d>0$, set
$A_N=\im(de_N(T))$ and $\pi:J_N\to A_N$. This standard construction is
explained in \cite{CaroccaRodriguez}. It gives
\begin{equation}\label{eq:rankone}
 H_1(A_N,\Q)\simeq\Q[X]/(\Phi_N(X)),\qquad
 \dim A_N=\varphi(N)/2.
\end{equation}
For every $1\le k<N/2$ with $\gcd(k,N)=1$, there is an algebraic
invariant form $\eta_k$ with $\pi^*\eta_k=\omega_k$. Explicitly, if
$i:A_N\hookrightarrow J_N$, take $\eta_k=d^{-1}i^*\omega_k$.
The rational action is irreducible, but $A_N$ need not be a simple abelian
variety. This distinction permits $N=8$ and composite exponents.

\begin{lemma}[Injective rational periods]\label[lemma]{lem:period-inject}
The period map
$\mathcal P:H_1(A_N,\Q)\to\C$, $\gamma\mapsto\int_\gamma\eta_1$,
is injective. If $U$ is a logarithm of an algebraic point of $A_N$ and
$\eta_1(U)=0$, then $U=0$.
\end{lemma}
\begin{proof}
As $\mathcal P(T_*\gamma)=\zeta_N\mathcal P(\gamma)$, its kernel is a
rational $T$-invariant subspace. Irreducibility of \eqref{eq:rankone}
makes the kernel zero or the whole space. In the latter case the nonzero
invariant form vanishes on the entire period lattice, which spans the real
Lie algebra, an impossibility.

Apply \cref{thm:AST} to $U$ and $\ker\eta_1$. A positive-dimensional
abelian subgroup there has a nonzero rational cycle, injective in the
homology of $A_N$, with zero period. The first assertion excludes this.
Thus the subgroup, and then $U$, is zero.
\end{proof}

\subsection{A second integral contradicts proportional division}
If $\varphi(N)\ge4$, there exists
$1<k<N/2$ with $\gcd(k,N)=1$. On the positive real path
$P(t)=(t,\sqrt{1-t^N})$, based at $O=(0,1)$, let $V(s)$ be the actual
Abel logarithm projected to $A_N$. Then
\[
 \exp_{A_N}V(s)=\pi([P(s)-O]),\qquad
 \eta_j(V(s))=I_{N,j}(s):=\int_0^s\frac{t^{j-1}\dd t}{\sqrt{1-t^N}}.
\]
All these integrals converge at $1$. If $q=a/b\in\Q\cap(0,1)$ and
$s=s_{N,q}$ were algebraic, the vector
$U=bV(s)-aV(1)$ would exponentiate to an algebraic point and would satisfy
$\eta_1(U)=0$. By \cref{lem:period-inject}, $U=0$.

But with $w_N(t)=(1-t^N)^{-1/2}$, the weighted mean
\[
 M_k(s)=\frac{I_{N,k}(s)}{I_N(s)}
\]
is strictly increasing, since
\begin{equation}\label{eq:weighted-monotone}
 M_k'(s)=\frac{w_N(s)}{I_N(s)^2}
 \int_0^s(s^{k-1}-t^{k-1})w_N(t)\dd t>0.
\end{equation}
It follows that
\[
 \eta_k(U)=bI_N(s)\{M_k(s)-M_k(1)\}<0,
\]
a contradiction. This proves the transcendental cases of
\cref{thm:all-division}.

\subsection{The small-exponent algebraic cases}
For $N=1$, $s_{1,q}=2q-q^2$. For $N=2$,
$s_{2,q}=\sin(\pi q/2)$, algebraic at rational $q$.
For $N=3,4$, the curve has genus one and the endpoint divisor $[W-O]$,
$W=(1,0)$, is torsion. For $N=3$, the point $O$ is 3-torsion relative to
infinity, since $y=1$ meets the cubic there with multiplicity three, while
$W$ is 2-torsion. For $N=4$, taking $O$ as origin, $(x,y)\mapsto(-x,y)$
is negation, $O'=(0,-1)$ is 2-torsion, and the hyperelliptic involution
implies $2W=O'$. Thus $W$ has order dividing four.

For $N=6$, the finite map
\[
 (x,y)\longmapsto(u,v)=(x^2,xy),\qquad v^2=u-u^4,
 \qquad du/v=2dx/y
\]
lands on a genus-one curve and sends both endpoints to branch points.
Their difference is 2-torsion. A rational multiple of a torsion logarithm
exponentiates to a torsion point. Such points on an algebraic elliptic curve,
and preimages under a finite algebraic map, are algebraic. These observations
prove all exceptional cases. They agree with the classical small-leaf
results reviewed in \cite{Tamiozzo}.

\subsection{The prime-only route and why it was superseded}
For a prime $p\ge5$, the curve $y^2=x-x^{p+1}$ is isomorphic over $\Qb$ to
$v^2=1-u^p$, by $u=1/x$, $v=\ii y/x^{(p+1)/2}$. Its Jacobian is absolutely
simple, as in \cite[Proposition 4.8]{Tamiozzo}. Along
$(x,y)=(s^2,s\sqrt{1-s^{2p}})$, the first two holomorphic forms pull back to
$2\dd s/\sqrt{1-s^{2p}}$ and
$2s^2\dd s/\sqrt{1-s^{2p}}$. Simplicity and the analytic subgroup theorem
then give the same weighted-mean contradiction. This was the prime-leaf
argument. The rational-period proof above removes the unnecessary
simplicity hypothesis and includes all composite exponents.

For $N=8$ the projector is particularly explicit:
\[
 \chi_8=(X^2+1)(X^4+1),\qquad e_8(T)=\frac{1-T^4}{2}.
\]
Its two-dimensional image uses $dx/y$ and $x^2dx/y$. This is compatible
with a Jacobian isogenous to a product of elliptic curves: irreducibility
of a rational rotation representation is not absolute simplicity.

\subsection{Geometric division points and their boundary exceptions}
The $n$-leaf curve $C_n=\{|z^n-1|=1\}$ has polar equation
$r^n=2\cos(n\theta)$ on the leaf about the positive real axis.
Its half-leaf arclength from the origin is
\begin{equation}\label{eq:leaf-length}
 \ell(r)=2^{1/n}I_{2n}(r/2^{1/n}).
\end{equation}
Thus every interior rational division point on a half-leaf is transcendental
in its normalized radius when $n\ge4$. Both Cartesian coordinates are
individually transcendental: if one were algebraic, the equation
\[
 (X^2+Y^2)^n-2\operatorname{Re}(X+\ii Y)^n=0
\]
is monic of degree $2n$ in the other coordinate, forcing that coordinate
and then the radius to be algebraic.

The origin and leaf tips remain algebraic. In a traversal starting at the
origin and visiting each leaf once, the $j$th point of an $M$-fold equal
arclength division lies at a half-leaf boundary exactly when
$M\mid2nj$. Hence, for $n\ge4$, all $M$-division points are algebraic
exactly when $M\mid2n$. Several boundary times can represent the same origin;
this is a statement about the specified traversal, not a count of distinct
plane points.

The first new normalized radii include
\begin{align*}
 s_{8,1/2}&=0.5813725649213314183423972021203471486707\ldots,\\
 s_{10,1/2}&=0.5660673118936379970848379116070460706217\ldots,\\
 s_{12,1/2}&=0.555521188888391003316451341835\ldots.
\end{align*}
For algebraic $0\le a<b\le1$, the same proof based at $P(a)$ shows that
rational division of $\int_a^b w_N(t)\dd t$ gives a transcendental interior
location whenever $\varphi(N)\ge4$.
Finally,
\[
 I_N(s)=\frac1N B_{s^N}(1/N,1/2),
\]
so the rational quantiles of $\operatorname{Beta}(1/N,1/2)$ are algebraic
exactly for $N=1,2,3,4,6$; in all other cases they are transcendental.

\section{Algebraic division ratios and torsion}\label{sec:ratios}
Rational division uses integer multiplication of logarithms. Algebraic
irrational multiplication is not automatically an algebraic-group operation.
The missing step is supplied by the period field.

\begin{lemma}[An algebraic ratio comes from an endomorphism]\label[lemma]{lem:ratio}
Let $A/\Qb$ have a finite-order automorphism $T$ whose characteristic
polynomial on rational homology is $\Phi_N$. Set
$K=\Q[T]\simeq\Q(\zeta_N)$. Let $\eta\ne0$ be an algebraic eigenform,
with eigenvalue corresponding to the embedding $\sigma:K\hookrightarrow\C$.
If $u,v$ are logarithms of algebraic points, $\eta(v)\ne0$, and
$\lambda=\eta(u)/\eta(v)$ is algebraic, then $\lambda\in\sigma(K)$ and
\begin{equation}\label{eq:log-endomorphism}
 u=da(v),\qquad a\in K,\quad\sigma(a)=\lambda.
\end{equation}
If $\exp_A(v)$ is torsion, algebraicity of this ratio is equivalent to
torsion of $\exp_A(u)$.
\end{lemma}
\begin{proof}
As in \cref{lem:period-inject}, the period map is injective. Since its domain
is a rank-one $K$-module, its image is $\Omega\sigma(K)$ for some
$\Omega\ne0$. Apply \cref{thm:AST} to $A\times A$, $(u,v)$, and the
hyperplane $\eta(Z_1)-\lambda\eta(Z_2)=0$. The resulting subgroup has a
positive-dimensional second projection, because $v\ne0$. Choose a rational
cycle with nonzero second projection. Its second period is nonzero by
injectivity, and integration of the hyperplane relation expresses $\lambda$
as the ratio of two elements of $\Omega\sigma(K)$. Thus
$\lambda\in\sigma(K)$.

Only now select the corresponding $a\in K\subseteq\End^0(A)$. Clearing
denominators, $Mu-d(Ma)(v)$ exponentiates to an algebraic point and is
annihilated by $\eta$. The last assertion of \cref{lem:period-inject} makes
it zero. This proves \eqref{eq:log-endomorphism}.
A torsion logarithm is a rational period vector, and quasi-endomorphisms
preserve rational periods. Conversely, the period-image description gives
an algebraic ratio of two torsion logarithms.
\end{proof}

\subsection{The reference endpoint is always torsion}
On $\mathcal H_N$, put $O=(0,1)$, $O'=(0,-1)$, and $W_r=(r,0)$ for
$r^N=1$. The divisor identities
\begin{equation}\label{eq:branch-torsion}
 \divisor\frac{y-1}{y+1}=N(O-O'),\qquad
 \divisor\frac{x-r}{x}=2W_r-O-O'
\end{equation}
show $2N[W_r-O]=0$. In particular $V(1)$ in \cref{sec:division} is a
nonzero logarithm of a torsion point after projection. For algebraic
$s\in(0,1)$, \cref{lem:ratio} therefore gives the exact equivalence
\begin{equation}\label{eq:ratio-torsion}
 R_N(s)\in\Qb
 \quad\Longleftrightarrow\quad
 \pi([P(s)-O])\text{ is torsion in }A_N.
\end{equation}
The right side concerns the projected point. Replacing it by torsion in
$J_N$ without an additional argument would be invalid.

\subsection{Complete exclusion for \texorpdfstring{$N=8,p,2p$}{N = 8, p, 2p}}
\begin{theorem}\label{thm:special-ratios}
Let $p\ge5$ be prime and $N\in\{8,p,2p\}$. Every algebraic
$s\in(0,1)$ has transcendental $R_N(s)$. Equivalently, every
algebraic-probability inverse value $s_{N,q}$ is transcendental.
\end{theorem}
\begin{proof}
For $N=8$, use $C:y^2=x-x^5$ with
$P(s)=(s^2,s\sqrt{1-s^8})$, $O=(0,0)$, $W=(1,0)$. Its genus-two
Jacobian has the automorphism $(x,y)\mapsto(\ii x,\zeta_8y)$ and full
homology characteristic polynomial $\Phi_8$. Moreover
$P^*(dx/y)=2\dd s/\sqrt{1-s^8}$ and $[W-O]$ is 2-torsion.
Thus an algebraic ratio forces $[P(s)-O]$ itself to be torsion.

Choose $\delta^2=\zeta_8$ and set $X=\zeta_8x$, $Y=\delta y$.
The equation becomes $Y^2=X^5+X$. Boxall--Grant's complete packet
\cite[Theorem (0.3)]{BoxallGrant} consists of infinity and the points whose
$X$-coordinates are roots of
\[
 X^5+X,\qquad X^4+4X^2+1,\qquad X^4-4X^2+1.
\]
There are 22 points in total, including infinity. The change of base point
from $(0,0)$ to infinity adds 2-torsion. On $x=s^2\in(0,1)$ the three
expressions become
\[
 \zeta_8x(1-x^4),\qquad 1-x^4+4\ii x^2,\qquad1-x^4-4\ii x^2,
\]
all nonzero. No point of our open arc belongs to the packet.

For $N=p$, the entire Jacobian of $y^2=1-x^p$ is the primitive factor.
With $X=-x$, $Y=y$ it becomes $Y^2=X^p+1$, and our arc has $-1<X<0$.
The base point differs from infinity by torsion because
$\divisor(y-1)=p(O-\infty)$. For $N=2p$, instead use
$y^2=x-x^{p+1}$ and the same squared-parameter path as above. Its full
homology is the $\Phi_{2p}$ module. The transformation
\[
 X=-1/x,\qquad Y=\ii y/x^{(p+1)/2}
\]
sends it to $Y^2=X^p+1$, sends $O$ to infinity, and sends our arc to $X<-1$.

Arul's classification \cite[Theorem 6.2.21]{Arul} says that for $p\ge7$
only infinity or a zero coordinate can give torsion relative to infinity.
For $p=5$ the additional orbit has $X^5=4$. None lies on either negative
real arc. This excludes torsion in both cases and proves the theorem.
\end{proof}

For example $R_8(1/2)=0.4297970681351117797271930073547706727140\ldots$
and $R_{10}(1/2)=0.4415953036626671314933075173404976638640\ldots$
are transcendental. This is stronger than their irrationality, which alone
would follow by contraposition from rational-division results.

\subsection{General algebraic exceptions: fields, conjugates, and finiteness}
Suppose $\varphi(N)\ge4$ and algebraic $s\in(0,1)$ has algebraic ratio
$q=R_N(s)$. The endomorphism in \cref{lem:ratio} lies in the real cyclotomic
field $K_N^+=\Q(\zeta_N+\zeta_N^{-1})$. Evaluating the logarithm identity
on each primitive holomorphic eigenform gives
\begin{equation}\label{eq:ratio-conjugates}
 \sigma_k(q)=\frac{I_{N,k}(s)}{I_{N,k}(1)},\qquad
 1\le k<N/2,\quad\gcd(k,N)=1.
\end{equation}
All these values lie in $(0,1)$. If $k<\ell$, then
\begin{align}
 &I_{N,k}(s)I_{N,\ell}(1)-I_{N,\ell}(s)I_{N,k}(1)\notag\\
 &\quad=\int_0^s\int_s^1
 (t^{k-1}r^{\ell-1}-t^{\ell-1}r^{k-1})w_N(t)w_N(r)\dd r\dd t>0.
\end{align}
Consequently the conjugates in \eqref{eq:ratio-conjugates} are distinct and
\begin{equation}\label{eq:ratio-degree}
 \Q(q)=K_N^+,\qquad [\Q(q):\Q]=\varphi(N)/2.
\end{equation}
Thus an algebraic probability with a conjugate outside $(0,1)$, which is
not totally real, or which has the wrong degree, cannot have an algebraic
inverse value. These are necessary, not sufficient, conditions for an
exception.

The exceptional set
\[
 \mathcal E_N=\{s\in\Qb\cap(0,1):R_N(s)\in\Qb\}
\]
is finite. Indeed, $Z_N=\pi([\mathcal H_N-O])$ is an irreducible image curve
containing zero and generating $A_N$. Its image is not a point because
$\pi^*\eta_1\ne0$. Since $\dim A_N\ge2$, it cannot be a torsion translate
of an elliptic subvariety. Manin--Mumford therefore gives finitely many
torsion points on $Z_N$ \cite{Roessler}. The map from the smooth projective
curve to its image is finite, and \eqref{eq:ratio-torsion} proves the claim.
No effective bound on the number of exceptions is claimed by this step.

\subsection{Two algebraic irrational ratios for every exponent}
For $q=(\sqrt5-1)/2$ or $q=1/\sqrt2$, a conjugate is negative. Thus
\eqref{eq:ratio-conjugates} excludes an algebraic inverse when
$\varphi(N)\ge4$. For $N=3,4,6$ the real endomorphism field is $\Q$, so
\cref{lem:ratio} likewise excludes either irrational ratio. For $N=2$ the
sine expression is excluded by Gelfond--Schneider. Hence both fixed ratios
have transcendental inverses for every $N\ge2$; $N=1$ is the explicit
algebraic exception $2q-q^2$.
For the golden-ratio division at $N=8$,
\begin{align*}
 s&=0.7163027194141435098452680809081946272424\ldots,\\
 s^8&=0.0693063136301807054228930664456009672487\ldots.
\end{align*}
The latter is the corresponding beta quantile.

\section{Local abelian logarithms and primitive divisors}\label{sec:local}
The torsion reformulation admits a different exclusion mechanism: exhibit a
single nonzero local logarithm. Complex and $\ell$-adic integral values are
not identified. It is the algebraic assertion of torsion that passes between
embeddings.

\begin{lemma}[Local nonvanishing, allowing ramification]\label[lemma]{lem:local}
Let $N\ge3$, let $\ell\ge5$ be prime with $\ell\nmid N$, and let an
algebraic $s\in(0,1)$ be a unit in a finite extension of $\Q_\ell$. Normalize
$v_\ell(\ell)=1$. If
\begin{equation}\label{eq:local-threshold}
 h=v_\ell(1-s^N)>\frac{2}{\ell-1},
\end{equation}
then $\pi([P(s)-O])$ is not torsion. In particular $R_N(s)$ is
transcendental.
\end{lemma}
\begin{proof}
Hensel lifting gives $r^N=1$ with $r\equiv s$ in the residue field. At
$W_r=(r,0)$ use $y$ as a parameter, and set $z=\sqrt{1-s^N}$. Then
$v_\ell(z)=h/2>1/(\ell-1)$ and
\[
 x=r(1-y^2)^{1/N},\qquad
 \frac{dx}{y}=-\frac{2r}{N}(1-y^2)^{1/N-1}\dd y.
\]
The curve has good reduction since $\ell\nmid2N$. In this small formal
neighborhood the abelian integral is its logarithm coordinate
\begin{equation}\label{eq:local-series}
 \mathcal L=-\frac{2r}{N}\sum_{j\ge0}
 \frac{(1-1/N)_j}{j!(2j+1)}z^{2j+1}.
\end{equation}
The functorial logarithm and its Newton--Leibnitz property are as in
\cite{Zarhin}; the displayed valuation places the formal group coordinates
inside the small logarithm convergence domain.

The coefficients $(1-1/N)_j/j!$ belong to $\Z_\ell$. If
$k=v_\ell(2j+1)$, then $1+k(\ell-1)\le\ell^k\le2j+1$, whence
\[
 v_\ell(2j+1)\le\frac{2j}{\ell-1}.
\]
The valuation of each term with $j\ge1$, minus that of the first term, is
at least
\[
 jh-v_\ell(2j+1)\ge j\left(h-\frac2{\ell-1}\right)>0,
\]
and this lower bound tends to infinity. The first term is uniquely minimal,
so $v_\ell(\mathcal L)=h/2$ and $\mathcal L\ne0$.

If $\pi([P(s)-O])$ were torsion, then so would
$\pi([P(s)-W_r])$, by \eqref{eq:branch-torsion}. Functoriality, torsion
annihilation, and $\pi^*\eta_1=dx/y$ would give
\[
 \mathcal L=\eta_1\bigl(\log_{A_N}\pi([P(s)-W_r])\bigr)=0,
\]
a contradiction. Now use \eqref{eq:ratio-torsion}.
\end{proof}
The strict inequality in \eqref{eq:local-threshold} is part of the statement.
If $s$ lies in an unramified extension and $v_\ell(1-s^N)\ge1$, it holds
automatically. This recovers the original unramified local test.

\subsection{All rational locations, without an exponent exception}
\begin{theorem}\label{thm:rational-input}
For every $N\ge3$ and every $s\in\Q\cap(0,1)$, $R_N(s)$ is
transcendental.
\end{theorem}
\begin{proof}
Write $s=a/b$, where $0<a<b$ and $\gcd(a,b)=1$. Bang--Zsigmondy provides
a primitive prime divisor $\ell$ of $b^N-a^N$, except at
$(b,a,N)=(2,1,6)$; this precise form is recalled in \cite{Avci}.
Primitivity makes the residue of $b/a$ have order $N$, so
$N\mid\ell-1$. Hence $\ell\ge5$, $\ell\nmid2Nab$, and
$v_\ell(1-s^N)\ge1$. Apply \cref{lem:local}.
At the sole exception, $2^6-1=63$ and $\ell=7$ still satisfies the local
conditions. A primitive divisor is a sufficient device for choosing the
prime, not a necessary hypothesis of the local lemma.
\end{proof}

For example, $R_{12}(1/2)$ can use $\ell=13$ because
$2^{12}-1=3^2\cdot5\cdot7\cdot13$. Its value is
\[
 0.4500171090998080361116676425739890811526\ldots.
\]
Other examples are $R_{18}(1/2)$, $R_{24}(1/2)$, and $R_{12}(2/3)$,
using primes $19$, $241$, and $61$, respectively. The geometric meaning is
fractional arclength at a specified rational normalized radius, not a
rational fraction of the arclength specified in advance.

\begin{corollary}
If $N=2^hM$ with $M\ge3$, and $s^{2^h}\in\Q\cap(0,1)$, then
$R_N(s)$ is transcendental.
\end{corollary}
\begin{proof}
Write $s^{2^h}=a/b$ and choose the good prime from the proof for $M$.
The polynomial $X^{2^h}-a/b$ has simple roots modulo this odd prime, and
all its roots lie in unramified extensions. Also
$1-s^N=1-(a/b)^M$ has positive integral valuation. Apply
\cref{lem:local}. The $M=6,a/b=1/2$ exception again uses $7$.
\end{proof}
For instance $R_{12}(2^{-1/4})=0.7606580391345716786952707914181839174367
\ldots$ is transcendental and equals the regularized incomplete beta value
$\mathrm I_{1/8}(1/12,1/2)$.

\section{Uniformity for bounded algebraic degree}\label{sec:bounded}
The preceding local lemma remains valid in ramified fields. Combining it
with a uniform primitive-ideal theorem gives a height-independent result.

\begin{external}[Voutier]\label[external]{thm:voutier}
Let $A,B$ be nonzero coprime algebraic integers in a number field, suppose
$A/B$ is not a root of unity, and put $d=[\Q(A/B):\Q]$. If
\[
 N>\max\{2(2^d-1),\ 4000(d\log(3d))^{12}\},
\]
then $B^N-A^N$ has a prime-ideal divisor which divides no earlier
$B^j-A^j$, $1\le j<N$.
\end{external}
This is \cite[Theorem 2(i)]{Voutier}, not the Lucas--Lehmer-only part of
that theorem. The degree in the bound is the degree of the ratio, rather
than the degree of an arbitrary ambient number field.

\begin{theorem}\label{thm:bounded}
For an integer $D\ge1$, define
\begin{equation}\label{eq:CD}
 C(D)=\max\{3^D,\ 4000(D\log(3D))^{12}\}.
\end{equation}
If algebraic $s\in(0,1)$ has degree at most $D$ and $N>C(D)$, then
$R_N(s)$ is transcendental. The bound is independent of height, denominators,
and the discriminant of $\Q(s)$.
\end{theorem}
\begin{proof}
Let $K=\Q(s)$, and write the fractional ideal $(s)=\mathfrak a\mathfrak
b^{-1}$ with coprime integral ideals. Some power
$\mathfrak b^h=(b_0)$ is principal. In a finite extension containing
$B^h=b_0$, one has $(B)=\mathfrak b\mathcal O_L$, because equality of
$h$th powers of fractional ideals implies equality. Set $A=sB$. Then
$A,B$ are coprime algebraic integers, while their ratio still has degree
$[K:\Q]$. Apply \cref{thm:voutier} in $L$.

Contract a primitive prime ideal to a prime $\mathfrak p$ of $K$. The
injection of residue fields preserves the exact order $N$ of $s$. If the
residue characteristic is $\ell$ and the residue degree in $K$ is $f$,
then
\[
 N\mid\ell^f-1,\quad \ell\nmid N,
 \quad v_{\mathfrak p}(1-s^N)\ge1.
\]
Let $e$ be its ramification index in the original field $K$. Then
$ef\le[K:\Q]\le D$ and $v_\ell(1-s^N)\ge1/e$.
If $\ell\le2e+1$, then
\[
 N+1\le\ell^f\le(2e+1)^f\le3^{ef}\le3^D,
\]
contradicting $N>C(D)$. Thus $\ell>2e+1$, so
$1/e>2/(\ell-1)$ and $\ell\ge5$. Apply \cref{lem:local}.
The comparison uses $e,f$ in $K$, not in the potentially much larger
principalizing extension.
\end{proof}

For $D=2$, $C(2)<20\,000\,000\,000$. This is explicit but very large;
no optimality is claimed. For each fixed algebraic $s$, only finitely many
exponents can give an algebraic ratio. More uniformly, an exceptional
sequence with exponents tending to infinity, if one exists, must have
input degrees tending to infinity.

\subsection{A small-exponent ramified quadratic example}
Let $s_0=3-\sqrt7\in(0,1)$. No positive power of $s_0$ is rational, since
its conjugate $3+\sqrt7>1$ has a different positive power. Thus it is not
covered by the power-root corollary above.
In $\Q_7(\sqrt7)$, $v_7(\sqrt7)=1/2$. If $N=6m$ with $7\nmid m$,
then $3^N\equiv1\pmod7$ and expansion gives
\[
 1-(3-\sqrt7)^N=1-3^N+N3^{N-1}\sqrt7+
 \sum_{j=2}^N(-1)^{j+1}\binom Nj3^{N-j}(\sqrt7)^j.
\]
The linear term has valuation $1/2$, whereas every other term has valuation
at least one. Hence $v_7(1-s_0^N)=1/2>1/3$, and
\begin{equation}
 R_{6m}(3-\sqrt7)\notin\Qb\qquad(m\ge1,\ 7\nmid m).
\end{equation}
In particular
\[
 1-(3-\sqrt7)^{12}=-524361663+198190080\sqrt7,
\]
and
\[
 R_{12}(3-\sqrt7)=0.3188329955550076211096403254230216298165\ldots.
\]
Here the local logarithm has valuation $1/4$. This small-exponent result
does not use the large constant \eqref{eq:CD}.

\section{Fermat beta integrals, polynomial weights, and second-kind forms}\label{sec:fermat}
The preceding sections exploit a particular hyperelliptic character. We now
change the curve rather than merely enlarge an exponent. A complete torsion
classification for Fermat quotients makes it possible to treat all algebraic
interior inputs for another family. The polynomial-weight argument then
requires a stronger statement than transcendence of a ratio: it requires
linear independence with an algebraic constant term.

\begin{theorem}[Beta independence and weighted stability]\label{thm:fermat-beta}
Let $p,r$ be distinct primes with $\max(p,r)\ge5$, and put
\[
 a=i/p,\quad b=j/r,\qquad 1\le i<p,\quad1\le j<r.
\]
No restriction $a+b<1$ is imposed. For every $x\in\Qb\cap(0,1)$,
\begin{equation}\label{eq:beta-three-independent}
 1,\quad \B(a,b),\quad \B_x(a,b)
 \quad\text{are linearly independent over }\Qb.
\end{equation}
Let $W(t)=\sum w_kt^k\in\Qb[t]$, and define
\begin{equation}\label{eq:weight-coefficient}
 c_W=\sum_k w_k\frac{(a)_k}{(a+b)_k}.
\end{equation}
Then, at every such algebraic $x$,
\begin{equation}\label{eq:weight-algebraic-iff}
 \int_0^xW(t)t^{a-1}(1-t)^{b-1}\dd t\in\Qb
 \quad\Longleftrightarrow\quad c_W=0.
\end{equation}
If $c_W\ne0$, the integral normalized by its value at $1$ is
transcendental. In particular, if $W$ has real algebraic coefficients and
is strictly positive on $(0,1)$, every interior algebraic-probability
quantile of its normalized density is transcendental.
\end{theorem}
All nonnegative integral parameter shifts are included by taking
$W=t^n(1-t)^m$. There is no bound on the degree of $W$ or on the degree or
height of $x$. Equal prime denominators, general composite denominators,
and the prime pair $\{2,3\}$ are not included.

\subsection{The curve and the pole that cannot be ignored}
Work on the smooth projective model of
\begin{equation}
 C:\quad U^p+V^r=1.
\end{equation}
It has a unique point at infinity and genus $(p-1)(r-1)/2$. The differential
\begin{equation}\label{eq:fermat-differential}
 \eta=U^{i-1}V^{j-r}\dd U
\end{equation}
is regular at every finite point. At $V=0$, differentiation of the curve
equation supplies the factor $V^{r-1}$ needed for regularity. At infinity,
\begin{equation}\label{eq:fermat-order}
 \ord_\infty\eta=pr-ri-pj-1.
\end{equation}
Thus $a+b<1$ gives a holomorphic form. If $a+b>1$, the form has a pole of
order at least two, with residue zero. Equality $a+b=1$ is impossible for
the specified prime denominators. The vanishing residue also follows
from the global residue theorem, since infinity is the only possible
pole. The real path
\[
 P_x=(x^{1/p},(1-x)^{1/r}),\qquad 0\le x\le1,
\]
satisfies
\begin{equation}\label{eq:fermat-pullback}
 P^*\eta=\frac1p x^{a-1}(1-x)^{b-1}\dd x.
\end{equation}
For example $a=2/3$, $b=2/5$ gives $C:U^3+V^5=1$ and
$\eta=UV^{-3}\dd U$, with $\ord_\infty\eta=-2$. Treating it as a
holomorphic differential on the ordinary Jacobian would be incorrect.

\subsection{Generalized Jacobians retain the polar part}
Choose $M\ge2$ at least as large as the pole order. The generalized
Jacobian for the modulus $M\infty$ fits into
\begin{equation}\label{eq:generalized-extension}
 0\longrightarrow \Ga^{M-1}\longrightarrow G
 \xrightarrow{\varrho}J\longrightarrow0,
 \qquad J=\Jac(C).
\end{equation}
There is no torus because the modulus has a single support point.
This is the standard modulus construction \cite{Zannier}. With
$O=(0,1)$ as base point, an algebraic Abel map
$j_O:C\setminus\{\infty\}\to G$ and an algebraic invariant differential
$\widetilde\eta$ satisfy
\begin{equation}
 j_O^*\widetilde\eta=\eta.
\end{equation}
One can check the differential condition locally rather than infer it
solely from an analytic formula. The pinched local ring is
$\Qb+t^M\Qb[[t]]$. The Rosenlicht condition asks that
$\operatorname{Res}(f\eta)=0$ for every element of this ring. Constants
are handled by the zero residue, and multiplication by $t^M$ removes the
pole. The corresponding dualizing differential therefore gives an
invariant differential of the generalized Jacobian. See \cite{ACL} for
the residue formulation and integral description of the Abel map.
All the modulus and Picard constructions used here are over $\Qb$.

The vector-group fiber in \eqref{eq:generalized-extension} is contractible
as a complex manifold. Consequently
\begin{equation}
 H_1(G(\C),\Q)\simeq H_1(J(\C),\Q).
\end{equation}
The rotation $T(U,V)=(\zeta_pU,\zeta_rV)$ preserves the modulus and acts on
$G$. Its action on the Abel map can differ by translation because it need
not fix $O$; invariant differentials are unaffected by that translation.
On $H_1(J,\Q)$ its eigenvalues are exactly
$\zeta_p^u\zeta_r^v$, $1\le u<p$, $1\le v<r$, each once. Therefore
\begin{equation}\label{eq:fermat-homology}
 \chi_{T_*}=\Phi_{pr},\qquad H_1(J,\Q)\simeq K=\Q(\zeta_{pr})
\end{equation}
as a rank-one $K$-module. This also follows from the homology calculation
in \cite[Proposition 2.4.4 and Corollary 2.4.5]{Arul}. Unlike a general
composite hyperelliptic exponent, the entire Jacobian is the relevant
rational character module.

\subsection{A nonzero period for a possibly second-kind form}
It is not enough here to say that the differential is nonzero: a nonzero
exact differential can have every closed period zero. We exhibit a
specific nonzero period.
Let $\ell$ be the real path from $O=(0,1)$ to $W=(1,0)$, and set
\[
 S(U,V)=(\zeta_pU,V),\qquad R(U,V)=(U,\zeta_rV).
\]
The map $S$ fixes $O$, and $R$ fixes both $W$ and $SW$. Hence
\[
 \Gamma=(1-R_*)(1-S_*)\ell
\]
is a closed integral cycle avoiding infinity. Directly,
\begin{equation}\label{eq:fermat-nonzero-period}
 \int_\Gamma\eta
 =(1-\zeta_r^j)(1-\zeta_p^i)\frac{\B(a,b)}p\ne0.
\end{equation}
The Abel map realizes this as a period on $G$. If $\sigma:K\hookrightarrow
\C$ sends $T$ to $\xi=\zeta_p^i\zeta_r^j$, the eigencharacter relation and
\eqref{eq:fermat-homology} now imply
\begin{equation}\label{eq:generalized-period-map}
 \ker\mathcal P_{\widetilde\eta}=0,
 \qquad\im\mathcal P_{\widetilde\eta}=\Omega\sigma(K),\quad\Omega\ne0.
\end{equation}
As before, this is injectivity on rational homology, not on its complex
scalar extension.

\subsection{Projection rigidity, not vanishing of the entire logarithm}
The additive directions in $G$ require a modification of
\cref{lem:ratio}. A logarithm can have algebraic coordinates in those
directions without being zero.

\begin{lemma}[Generalized projection rigidity]\label[lemma]{lem:projection}
If $w\in\Lie(G)_\C$ exponentiates to an algebraic point and
$\widetilde\eta(w)\in\Qb$, then $\dd\varrho(w)=0$.
If $u,v$ exponentiate to algebraic points, $\bar v=\dd\varrho(v)\ne0$, and
\begin{equation}\label{eq:affine-two-log}
 \widetilde\eta(u)-\lambda\widetilde\eta(v)=c,
 \qquad\lambda,c\in\Qb,
\end{equation}
then $\lambda\in\sigma(K)$ and
\begin{equation}\label{eq:projection-relation}
 \dd\varrho(u)=a_*\bar v
\end{equation}
for the corresponding $a\in K\subseteq\End^0(J)$.
\end{lemma}
\begin{proof}
For the first assertion, let $c=\widetilde\eta(w)$ and apply
\cref{thm:AST} to $(w,1)$ on $G\times\Ga$ and the hyperplane
$\widetilde\eta(Z)-cT_0=0$. Let $H$ be the resulting connected subgroup.
If its image $B\subseteq J$ were positive dimensional, the kernel of
$H\to B$ would be a vector group. Hence rational homology maps onto
$H_1(B,\Q)$, and a cycle with nonzero image in $J$ could be chosen.
The integral of the additive coordinate differential around this cycle
is zero. The hyperplane then forces its $\widetilde\eta$-period to vanish,
contrary to \eqref{eq:generalized-period-map}. Thus $B=0$.

For the second assertion, use $G\times G\times\Ga$, the logarithm
$(u,v,1)$, and \eqref{eq:affine-two-log}. The second abelian projection
is positive dimensional since $\bar v\ne0$. Choose a rational cycle
whose second abelian projection is nonzero. Integrating the hyperplane
relation gives $\lambda$ as a ratio of two periods in
$\Omega\sigma(K)$, with nonzero denominator. Thus $\lambda\in\sigma(K)$.

It is not necessary to embed the field $K$ into the entire endomorphism
algebra of $G$. Choose $A(X)\in\Q[X]$ with $A(\xi)=\lambda$ and use the
actual quasi-endomorphism $A(T)$ of $G$. It induces the required $a$ on
$J$. After clearing denominators, $mu-(mA(T))_*v$ exponentiates to an
algebraic point and has algebraic $\widetilde\eta$-coordinate $mc$.
The first assertion gives \eqref{eq:projection-relation}.
\end{proof}
If $\bar v$ is a torsion logarithm, it is a rational period, and
\eqref{eq:projection-relation} makes $\dd\varrho(u)$ a rational period
as well. Only the image in the ordinary Jacobian is forced to be torsion;
no such claim is made about the point of $G$ itself.

\subsection{Torsion packets force three-term independence}
The divisors
\begin{equation}
 \divisor(V-1)=p(O-\infty),\qquad
 \divisor(U-1)=r(W-\infty)
\end{equation}
show that $[W-O]$ is torsion in $J$. Choose actual logarithms $u_x,v=u_1$
along the real path. Then
\begin{equation}
 \widetilde\eta(u_x)=\B_x(a,b)/p,
 \qquad\widetilde\eta(v)=\B(a,b)/p.
\end{equation}
The projected logarithm $\bar v$ is a nonzero rational period. Its
nonvanishing is independently witnessed by the positive integral of the
holomorphic differential with indices $i=j=1$, since $1/p+1/r<1$.

Suppose $c_0+c_1\B(a,b)+c_2\B_x(a,b)=0$ with algebraic coefficients.
If $c_2\ne0$, \cref{lem:projection} makes $[P_x-O]$ torsion in $J$,
and hence also $[P_x-\infty]$ is torsion. For odd $p,r$, the transformation
$(X,Y)=(-U,V)$ gives $Y^r=X^p+1$. Arul's complete classification
\cite[Theorem 6.2.21]{Arul} excludes noncuspidal torsion points for
these distinct odd prime exponents. Our point has both coordinates
nonzero, so it is excluded.

If one prime is $2$, interchange variables if necessary and write
$r=2$, $p\ge5$. The real image has $X=-x^{1/p}\in(-1,0)$. For
$p\ge7$ there is no additional torsion orbit. For $p=5$ the additional
orbit satisfies $X^5=4$, and therefore misses this interval. These are
precisely the exceptional cases retained from the classification, not
ignored by it. Thus $c_2=0$. If $c_1\ne0$, the beta period would be
algebraic, and the first part of \cref{lem:projection} applied to $v$
would contradict $\bar v\ne0$. Hence all three coefficients vanish,
proving \eqref{eq:beta-three-independent}.

In the holomorphic subcase the same reasoning can be performed directly
on $J$: \cref{lem:ratio} supplies the torsion reduction and
\cref{lem:affine-separation} removes the algebraic constant term.
The generalized argument is what permits $a+b>1$.

\subsection{Exact polynomial reduction and its algebraic exception}
Set $D(x)=x^a(1-x)^b$. Standard incomplete-beta recurrences
\cite{DLMF} give a polynomial $P_W\in\Qb[x]$ such that
\begin{equation}\label{eq:weighted-beta-reduction}
 \int_0^xW(t)t^{a-1}(1-t)^{b-1}\dd t
 =c_W\B_x(a,b)+D(x)P_W(x).
\end{equation}
For clarity, a constructive recurrence for $W=t^k$ is
\begin{align}
 c_0&=1,&P_0&=0,\
 c_{k+1}&=\frac{a+k}{a+b+k}c_k,
 &P_{k+1}&=\frac{a+k}{a+b+k}P_k-\frac{x^k}{a+b+k}.
\end{align}
Linear combinations handle every polynomial. Differentiation verifies
\begin{equation}\label{eq:weight-differential}
 x(1-x)P_W'+[a-(a+b)x]P_W=W-c_W.
\end{equation}
The endpoint term vanishes at both $0$ and $1$, so the total integral is
$c_W\B(a,b)$. At algebraic $x$, $D(x)P_W(x)$ is algebraic. Thus
\eqref{eq:beta-three-independent} proves
\eqref{eq:weight-algebraic-iff}; it also excludes an algebraic normalized
value when $c_W\ne0$. Positivity gives $c_W>0$ and strict monotonicity,
which supplies the quantile statement. This completes the proof of
\cref{thm:fermat-beta}.

For a fixed degree bound on $W$, the algebraic cases form exactly the
linear hyperplane $c_W=0$. For example, with $a=2/3,b=2/5$,
\begin{equation}\label{eq:signed-beta-exact}
 \int_0^x(8t-5)t^{-1/3}(1-t)^{-3/5}\dd t
 =-\frac{15}{2}x^{2/3}(1-x)^{2/5}.
\end{equation}
The weight changes sign and is not a probability density. The example
prevents an incorrect claim that every nonzero weighted integral is
transcendental.

\subsection{The exact linear dimension of all weighted values}
For fixed algebraic $x\in(0,1)$ and admissible positive $W$, write
\[
 F_0(x)=\B_x(a,b)/\B(a,b),\qquad
 F_W(x)=F_0(x)+\frac{h_W(x)}{\B(a,b)},\quad h_W(x)\in\Qb.
\]
Consequently
\begin{equation}\label{eq:weight-span}
 \dim_{\Qb}\Span\bigl(\{1\}\cup\{F_W(x):W>0\}\bigr)=3.
\end{equation}
The upper bound follows from the three displayed directions. Their
independence follows from \eqref{eq:beta-three-independent} after
multiplication by $\B(a,b)$. Taking $W=1,t$ realizes the third direction,
since
\[
 F_t(x)=F_0(x)-\frac{D(x)}{a\B(a,b)}.
\]
Equation \eqref{eq:weight-span} concerns linear dimension, not
transcendence degree.

\subsection{Explicit first-kind and second-kind examples}
For $a=1/3,b=1/5$, let $B_0=\B(a,b)$ and $D=x^{1/3}(1-x)^{1/5}$.
If $G$ denotes the $\operatorname{Beta}(4/3,6/5)$ distribution function,
and $H$ the distribution obtained with weight $1+t^2$, then
\begin{align}
 B_0G(x)&=\B_x(1/3,1/5)+(8x-3)D(x),\
 B_0H(x)&=\B_x(1/3,1/5)-\frac{15}{71}(5+2x)D(x).
\end{align}
Here $c_{t(1-t)}=15/184$ and $c_{1+t^2}=71/46$. In particular
\[
 71H(1/2)-161F_0(1/2)+90G(1/2)=0,
\]
whereas $1,F_0(1/2),G(1/2)$ are linearly independent. The median of $G$ is
\[
 0.5341541616046073020427688226978928570926\ldots,
\]
and is transcendental by \cref{thm:fermat-beta}.

For the second-kind example $a=2/3,b=2/5$, set
$B_0=\B(2/3,2/5)$ and $D=x^{2/3}(1-x)^{2/5}$. The distribution function
$G=\mathrm I_x(5/3,7/5)$ satisfies
\begin{equation}\label{eq:second-kind-median}
 B_0G(x)=\B_x(2/3,2/5)+(4x-3/2)D(x),
\end{equation}
with $c_{t(1-t)}=15/124$. Its median is
\begin{equation}
 m=0.5539892840285170056182880743334986529370\ldots.
\end{equation}
If $m$ were algebraic, \eqref{eq:second-kind-median} at $G(m)=1/2$
would contradict \eqref{eq:beta-three-independent}. This is not covered
by the integral-parameter-sum HPD theorem: $5/3+7/5=46/15$.

The hypotheses remain material. Equal denominators allow symmetric laws,
for example $\operatorname{Beta}(1/3,1/3)$ with median $1/2$.
The pair $\{2,3\}$ leads to an elliptic curve with noncuspidal torsion
points and algebraic rational divisions. No classification for arbitrary
rational-function weights or general composite denominators is asserted.
Dhillon--Saha \cite{DhillonSaha} is related work on incomplete-beta values;
a complete comparison with every theorem of that paper is not part of
this consolidation and no priority conclusion is drawn from its abstract.

\section{Attempted transfer to the Euler--Mascheroni constant}\label{sec:gamma}
This section records an attempted transfer, not a proof of the irrationality
or transcendence of $\gamma$. It also does not certify every preceding
manuscript argument by virtue of having compared it with the gamma problem.
The purpose is to determine precisely which hypotheses of the successful
mechanisms survive and which additional arithmetic implication is missing.

The target is the constant in \eqref{eq:gamma-def}, rather than the
function $\Gamma(z)$ as an element of a differential field. For example,
$\Gamma(n)=(n-1)!$ is algebraic at every positive integer, despite the
transcendental nature of the function. Likewise, transcendence of some
higher derivatives does not identify the arithmetic of $\Gamma'(1)$.

\subsection{A normalization ambiguity invisible to beta quotients}
The standard connection is
\begin{equation}
 \B(a,b)=\frac{\Gamma(a)\Gamma(b)}{\Gamma(a+b)}.
\end{equation}
For an arbitrary $c\in\C$, put $\Gamma_c(z)=e^{cz}\Gamma(z)$. Then
\begin{equation}\label{eq:beta-gauge}
 \frac{\Gamma_c(a)\Gamma_c(b)}{\Gamma_c(a+b)}=\B(a,b),
 \qquad \psi_c(z)=c+\psi(z),\quad\psi_c(1)=c-\gamma.
\end{equation}
Thus the beta quotient alone cannot detect the linear term in the
logarithm of the gamma function. In particular,
\[
 \partial_a\log\B(a,b)=\psi(a)-\psi(a+b)
\]
cancels the common constant $-\gamma$.
This is an information limitation of a beta-quotient-only transfer, not a
claim that $\Gamma_c$ satisfies every normalization or recurrence of
$\Gamma$. In general it does not. Adding normalization or growth data
changes the problem; \eqref{eq:beta-gauge} does not rule out every
possible argument involving beta integrals.

\subsection{Confluence returns gamma, but not an arithmetic lower bound}
For $\Re z>0$, define
\begin{equation}
 G_M(z)=M^z\B(z,M+1)
 =\int_0^M t^{z-1}(1-t/M)^M\dd t
 =\frac{M^zM!}{z(z+1)\cdots(z+M)}.
\end{equation}
Dominated convergence gives $G_M(z)\to\Gamma(z)$. The Laurent expansion
at zero is exact to first order:
\begin{equation}
 G_M(z)=\frac1z+\log M-H_M+O(z).
\end{equation}
Its finite part tends to $-\gamma$. This is a valid analytic connection,
but it returns to the classical harmonic-number approximation. Although
$H_M-\log M$ is transcendental for $M>1$, that fact alone does not classify
its limit. The same warning applies to the regularized $\Ein$ limit in
\eqref{eq:ein-gamma-limit}: its explicit appearance of $\gamma$ and a
dilogarithm is an identity, not an automatic transcendence transfer.

\subsection{Polynomial logarithmic moments lock the gamma and tail coefficients}
The polynomial-weight method can be tested directly. Let
$P(t)=\sum_{k=0}^d p_kt^k\in\Qb[t]$ and set
\begin{equation}
 B_P=\sum_{k=0}^d p_k k!,\qquad
 A_P=\sum_{k=0}^d p_k k!H_k.
\end{equation}
Differentiating the gamma recurrence gives
$\Gamma'(k+1)=k!(H_k-\gamma)$ \cite{DLMF}. Therefore
\begin{equation}\label{eq:gamma-moment}
 J_P:=\int_0^\infty P(t)e^{-t}\log t\dd t=A_P-B_P\gamma.
\end{equation}
These integrals converge absolutely. If $B_P=0$, $J_P$ is algebraic. If
$B_P\ne0$, then
\begin{equation}\label{eq:moment-equivalence}
 J_P\in\Qb\quad\Longleftrightarrow\quad\gamma\in\Qb.
\end{equation}
In particular, increasing the complexity of a polynomial does not by itself
produce an independently established transcendence result for such a
moment. For a positive real polynomial weight,
$B_P=\int_0^\infty P(t)e^{-t}\dd t>0$, so the question remains exactly
as strong as the gamma question.

To identify the tail, define the polynomials
\begin{equation}
 Q_P(t)=\sum_{j=0}^dP^{(j)}(t),\qquad
 R_P(t)=\frac{Q_P(t)-Q_P(0)}t,
 \qquad C_P=\sum_{j\ge0}R_P^{(j)}(1).
\end{equation}
The last sum is finite; it is zero when $R_P=0$. Directly,
\[
 Q_P-Q_P'=P,\qquad Q_P(0)=B_P,\qquad
 \int_0^\infty e^{-t}R_P(t)\dd t=A_P.
\]
The last identity follows first for monomials and then by linearity.
Integration by parts, using $\log1=0$, gives
\begin{align}
 T_P&:=\int_1^\infty P(t)e^{-t}\log t\dd t
   =\int_1^\infty e^{-t}\frac{Q_P(t)}t\dd t\notag\\
   &=B_PE_1(1)+\frac{C_P}{e}.\label{eq:gamma-tail}
\end{align}
With $\eta=\Ein(1)=\gamma+E_1(1)$, the lower piece is
\begin{equation}\label{eq:gamma-lower}
 L_P:=-\int_0^1P(t)e^{-t}\log t\dd t
 =B_P\eta-A_P+\frac{C_P}{e}.
\end{equation}
The signs agree with $J_P=T_P-L_P$.

\begin{center}
\begin{tabular}{@{}ccccc@{}}
\toprule
$P(t)$ & $A_P$ & $B_P$ & $C_P$ & $J_P$\\
\midrule
$1$ & $0$ & $1$ & $0$ & $-\gamma$\\
$t$ & $1$ & $1$ & $1$ & $1-\gamma$\\
$t^2$ & $3$ & $2$ & $4$ & $3-2\gamma$\\
$t^3$ & $11$ & $6$ & $17$ & $11-6\gamma$\\
$t-1$ & $1$ & $0$ & $1$ & $1$\\
\bottomrule
\end{tabular}
\end{center}
The same coefficient $B_P$ multiplies $-\gamma$ in
\eqref{eq:gamma-moment} and $E_1(1)$ in \eqref{eq:gamma-tail}.
Consequently any finite linear combination cancelling the displayed
$E_1(1)$ coefficient also cancels the gamma coefficient. This statement
is at the level of the explicit coefficients; it does not assume an
unproved algebraic relation or independence between $E_1(1)$ and $e$.
An exact counterexample to indiscriminate transcendence claims is
\begin{equation}\label{eq:gamma-algebraic-counterexample}
 \int_0^\infty(t-1)e^{-t}\log t\dd t=1.
\end{equation}
The scope is fixed $e^{-t}$, one factor $\log t$, finite polynomial
weights, and linear elimination. Rational-function weights, nonlinear
combinations, several logarithms, and new height estimates are not
excluded by this calculation.

\subsection{Why a finite modulus does not directly accommodate the exponential}
In \cref{sec:fermat}, the pole order was finite and a modulus $M\infty$
produced a finite-dimensional generalized Jacobian. In
\[
 -\gamma=\int_0^\infty e^{-t}\log t\dd t,
\]
the substitution $u=1/t$ introduces $e^{-1/u}$, an essential singularity.
A finite algebraic covering does not turn that exponential essential
singularity into a finite-order pole. Thus the current expression is
not made into an algebraic meromorphic differential by the earlier
modulus construction.
This is not a proof that $\gamma$ has no other period representation.
It identifies a failure of this particular direct construction. Moreover,
assuming $\gamma$ algebraic does not automatically supply an algebraic
point whose logarithm is the displayed integral, nor a torsion point on
the Fermat curves. Those were essential hypotheses of the preceding
analytic-subgroup and local arguments.

\subsection{Two arithmetic blocks containing the target exactly}
For positive real $z$,
\begin{equation}
 \Ein(z)=\gamma+\log z+E_1(z).
\end{equation}
Define the entire $E$-function
\begin{equation}\label{eq:gamma-A}
 A(z)=e^z\Ein(z)=\sum_{n\ge0}\frac{H_n}{n!}z^n
\end{equation}
and the formal factorial series at infinity
\begin{equation}\label{eq:gamma-D}
 \widehat D(z)=\sum_{n\ge0}(-1)^n n!z^{-n-1}.
\end{equation}
The latter diverges for every fixed nonzero complex $z$, but its
positive-direction Borel--Laplace sum is
\begin{equation}\label{eq:Borel-D}
 D_0(z)=\int_0^\infty\frac{e^{-t}}{z+t}\dd t=e^zE_1(z).
\end{equation}
For each algebraic candidate $r$, form
\begin{equation}
 \widehat\Psi_r(z)=A(z)-r e^z-\widehat D(z).
\end{equation}
Its defining coefficients use harmonic numbers, factorials, and $r$;
they do not require a decimal approximation to $\gamma$.
Positive-direction summation gives the exact formula
\begin{equation}\label{eq:gamma-target}
 \Psi_{r,0}(z)=e^z(\gamma-r+\log z),\qquad
 \Psi_{r,0}(1)=e(\gamma-r).
\end{equation}
The arithmetic framework for such mixed objects is that of
Fischler--Rivoal \cite{FischlerRivoalGevrey}. The formal mixed class
must be distinguished from an arbitrary analytic solution of a
differential equation.

\subsubsection{A fixed differential equation is not a value theorem}
The identities
\[
 A'=A+(e^z-1)/z,\qquad D_0'=D_0-1/z
\]
give $\Psi_{r,0}'=\Psi_{r,0}+e^z/z$. Eliminating $e^z$ yields
\begin{equation}\label{eq:gamma-ode}
 zY''+(1-2z)Y'+(z-1)Y=0.
\end{equation}
Every algebraic candidate $r$ gives a solution of the same equation,
whose solution space is spanned by $e^z$ and $e^z\log z$.
The coefficient $\gamma-r$ is a connection constant specifying the
arithmetic solution, not a coefficient of \eqref{eq:gamma-ode}.
Finding this equation therefore does not by itself exclude any candidate
$r$. The unresolved issue concerns compatibility of the two arithmetic
blocks, not existence of the analytic solution.

\subsection{An arithmetic division obstruction and the missing implication}
There is a concrete obstruction on one side of a potential zero theorem.

\begin{proposition}\label[proposition]{prop:gamma-division}
For algebraic $r$, there is no identity
\begin{equation}\label{eq:forbidden-division}
 \widehat\Psi_r(z)=(z-1)
 \left(F(z)+\sum_{n\ge1}b_nz^{-n}\right)
\end{equation}
with $F$ an $E$-function and all $b_n$ algebraic. In particular such a
factorization does not exist in the mixed arithmetic class.
\end{proposition}
\begin{proof}
Comparison of the nonnegative powers gives
\[
 A(z)-r e^z=(z-1)F(z)+b_1.
\]
Evaluating this entire-function identity at $1$ gives
$b_1=e\Ein(1)-re$. This is incompatible with the algebraic independence
of $e$ and $\Ein(1)$, a subcase of \cref{thm:ein} (or the classical
$E$-function argument in \cite{RivoalSurvey}).
\end{proof}
The missing implication is
\begin{equation}\label{eq:missing-zero-lift}
 \Psi_{r,0}(1)=0
 \quad\Longrightarrow\quad
 \widehat\Psi_r=(z-1)\times\text{a mixed arithmetic function}.
\end{equation}
An analytic zero permits analytic division on that branch. It does not
prove that the quotient has the required algebraic coefficients and
arithmetic denominator growth. A general zero-division assertion of this
kind is Conjecture 3 in \cite{FischlerRivoalGevrey}, not an established
input available for unconditional use.

In the present family, asserting \eqref{eq:missing-zero-lift} for every
algebraic $r$ is equivalent to transcendence of $\gamma$. One implication
uses \cref{prop:gamma-division} and \eqref{eq:gamma-target}; the other is
vacuous because if $\gamma$ is transcendental the premise never holds
for algebraic $r$. Thus this is a sharply formulated target, but not a
claim to have reduced the problem to an independently easier theorem.
The same obstruction had already appeared in the earlier exploratory
record recalled during the research conversation; it is not counted as
a newly discovered result in this consolidation.

There is also an instructive monodromy check. Under the hypothetical
identification $r=\gamma$, the quotient on the positive branch is
\[
 \frac{e^z\log z}{z-1},
\]
which has a removable singularity at $1$. Continuing once around $0$
adds $2\pi\ii e^z/(z-1)$, introducing a pole on the other branch.
This separates branchwise removability from a global arithmetic
regularity statement; it is not itself a transcendence proof.

\subsection{The two direct local series have disjoint convergence domains}
A local nonvanishing calculation succeeded in \cref{sec:local} because
an algebraic torsion assertion could be tested in one small residue disk.
The two most direct gamma series do not even converge at a common nonzero
local input.

\begin{proposition}\label[proposition]{prop:gamma-local}
Fix a prime $\ell$, an algebraic $\ell$-adic embedding of $r$, and normalize
$v_\ell(\ell)=1$. The series
\[
 A_r(z)=\sum_{n\ge0}\frac{H_n-r}{n!}z^n,
 \qquad\widehat D(z)=\sum_{n\ge0}(-1)^nn!z^{-n-1}
\]
converge, respectively, exactly on the open regions
\begin{equation}\label{eq:gamma-radii}
 v_\ell(z)>\frac1{\ell-1},
 \qquad
 v_\ell(z)<\frac1{\ell-1}.
\end{equation}
Neither converges on the common boundary.
\end{proposition}
\begin{proof}
Legendre's formula is
\[
 v_\ell(n!)=\frac{n-s_\ell(n)}{\ell-1},
\]
where $s_\ell(n)$ is the digit sum in base $\ell$. Also
$v_\ell(H_n)\ge-\lfloor\log_\ell n\rfloor$, so the harmonic coefficient
changes only a logarithmic correction to the factorial radius.
For necessity, take $n=\ell^k$. The unique smallest-valuation term in
$H_{\ell^k}$ is $1/\ell^k$, hence
\[
 v_\ell(H_{\ell^k})=-k,
 \qquad v_\ell(H_{\ell^k}-r)=-k
\]
for sufficiently large $k$. At the boundary
$v_\ell(z)=1/(\ell-1)$, the corresponding $A_r$ term has valuation
$-k+1/(\ell-1)$, whereas the corresponding $\widehat D$ term has
valuation $-2/(\ell-1)$. Neither sequence of terms tends to zero.
The same subsequence excludes the wrong side of each boundary, and the
factorial estimates give convergence on the stated sides.
\end{proof}
This does not prohibit $\ell$-adic analytic continuation, Frobenius
structures, or an overconvergent construction. It rules out a direct
comparison of these two defining series at the same nonzero local point.
Moreover a local nonvanishing statement alone would still require a
valid arithmetic bridge from the complex gamma hypothesis.

\subsection{A genuine conditional conclusion, kept conditional}
Let $\mathcal E$ be the ring of algebraic-point values of $E$-functions,
and let $\mathcal D$ be the ring of suitably summed arithmetic Gevrey
series of order $+1$, in the precise framework of
\cite{FischlerRivoalGevrey}. The following is an existing conjecture.
\begin{conjecture}[Fischler--Rivoal]\label[conjecture]{conj:ED}
\begin{equation}
 \mathcal E\cap\mathcal D=\Qb.
\end{equation}
\end{conjecture}
It implies the transcendence of $\gamma$. Indeed, put
\[
 \delta=eE_1(1)\in\mathcal D,
 \qquad\eta=\Ein(1),\qquad e,\eta\in\mathcal E.
\]
The identity
\begin{equation}\label{eq:gamma-delta}
 \delta=e\eta-e\gamma
\end{equation}
shows that an algebraic $\gamma$ would put $\delta$ in $\mathcal E$ as
well. Under \cref{conj:ED}, $\delta$ would then be algebraic, contradicting
the algebraic independence of $e$ and $\eta$. This is a specialization of
the more general conditional gamma-derivative result in
\cite[Theorem 2]{FischlerRivoalGevrey}, not a proof of the conjecture.

An unconditional but weaker consequence is already immediate from the
same algebraic independence: $\gamma$ and $\delta$ cannot both be
algebraic. This known alternative does not determine which of them is
transcendental \cite{RivoalSurvey}. Similarly,
\[
 \Gamma''(1)=\gamma^2+\frac{\pi^2}{6}
\]
is entirely compatible with algebraic $\gamma$ and transcendental
$\Gamma''(1)$. Even an infinitude of transcendental higher derivatives
would not by itself identify the first derivative.

\subsection{What survives the attempted transfer}
\begin{longtable}{@{}P{0.34\textwidth}P{0.60\textwidth}@{}}
\toprule
Mechanism examined & Outcome and exact limitation\\
\midrule
\endhead
Beta quotient and parameter differentiation & The common linear term of
$\log\Gamma$ is invisible to the quotient; normalization information must
be added.\\
Confluent beta limit & The finite part returns $H_M-\log M$; no new
arithmetic lower bound for its limit has been supplied.\\
Polynomial one-log moments & The same coefficient controls $\gamma$ and
the displayed exponential-integral tail; finite linear cancellation removes
both.\\
Finite-modulus generalized Jacobian & The exponential essential singularity
is not a finite-order algebraic pole in the given representation.\\
Mixed $E$/factorial construction & A fixed second-order equation contains
$e(\gamma-r)$ as the value of a selected arithmetic solution.\\
Arithmetic division & Formal division is excluded, but promotion of a
summed zero to arithmetic divisibility remains unproved.\\
Direct local evaluation & The two natural series have disjoint convergence
regions, including failure at the boundary.\\
Intersection conjecture & It would prove transcendence of $\gamma$, but
is not established here.\\
Arithmetic status of $\gamma$ & Neither irrationality nor transcendence is
proved by this manuscript.\\
\bottomrule
\end{longtable}
The missing work is an arithmetic noncancellation theorem for the selected
mixed solution, or a sufficiently strong quantitative substitute.
Repeating an unproved lifting assertion under a different name is not
progress toward that theorem. Nor does abstract nonsplitting of an
extension force a particular entry of an arbitrary comparison matrix to
be transcendental. These distinctions are essential when evaluating a
proposed transfer rather than merely producing another exact identity.

\section{Numerical certification, reproducibility, and limitations}\label{sec:verification}
Three logically distinct layers occur in this work: an external theorem,
a symbolic argument checking its hypotheses, and a finite computation.
A successful numerical test does not prove the external theorem or turn a
flawed hypothesis check into a valid proof. Conversely, an exact root
interval can be certified without assuming that its center is an accurate
floating-point approximation.

\subsection{Positive series and explicit remainder estimates}
For $0<a,b<1$ and $0<x\le1/2$,
\begin{equation}\label{eq:beta-positive-series}
 \B_x(a,b)=x^a\sum_{j\ge0}
 \frac{(1-b)_j}{j!(a+j)}x^j.
\end{equation}
Every term is positive and the ratio of consecutive terms is less than
$x$. If $T$ is the first omitted term, the remainder is at most
$T/(1-x)$. The complete integral is evaluated using
\[
 \B(a,b)=\B_{1/2}(a,b)+\B_{1/2}(b,a),
\]
and values with $x>1/2$ are obtained by subtraction from the complete
integral. Algebraic powers are enclosed by integer-root computations
on a rational grid. Polynomial corrections such as
\eqref{eq:weighted-beta-reduction} are then evaluated with rational interval
arithmetic, retaining the rounding error at every step.

Similarly,
\begin{equation}\label{eq:arc-cert-series}
 I_{N,k}(s)=\int_0^s\frac{t^{k-1}\dd t}{\sqrt{1-t^N}}
 =\sum_{j\ge0}\frac{\binom{2j}{j}s^{Nj+k}}{4^j(Nj+k)}
\end{equation}
has positive terms with consecutive ratio less than $s^N$ when $s<1$.
The tail is bounded by the first omitted term divided by $1-s^N$.
For the endpoint $1$, use
$I_{N,k}(1)=N^{-1}\B(k/N,1/2)$ and the preceding half-interval method.
No integration across an unbounded endpoint singularity is performed
numerically without a remainder bound.

For Gaussian primitives, Taylor's theorem gives
\begin{equation}
 \left|J_1(x)-\sum_{j=0}^m
 \frac{(-1)^jx^{2j+1}}{(2j+1)j!}\right|
 \le\frac{x^{2m+3}}{(2m+3)(m+1)!},\qquad x\ge0,
\end{equation}
with the sign of the first omitted term. Negative arguments are handled
by oddness. Arctangents are reduced to arguments of absolute value less
than one and enclosed by alternating series. A rational enclosure of
$\pi$ follows from
\begin{equation}
 \pi=16\arctan(1/5)-4\arctan(1/239).
\end{equation}
The lemniscate calculation uses the exact coefficient-square sum and
Cauchy bound in \eqref{eq:L3-series}--\eqref{eq:L3-tail}.

For a proposed root interval $[u,v]$, the verifier proves that the
residual is strictly negative at $u$ and strictly positive at $v$.
The mathematical monotonicity theorem supplies existence and uniqueness
inside the interval. The printed midpoint is never used as a proof of
its own correctness.

\subsection{Verification scripts rerun for the unified edition}
The fourteen verification programs contained in the supplied research
archives were inspected and rerun with their default exact-arithmetic
settings during preparation of this edition. All fourteen completed
successfully. A second execution of each program produced a byte-identical
JSON result, checked by SHA-256. The optional floating-point mode of the
gamma-transfer program was not needed for these reruns.

The following table identifies the checks and records short prefixes of
the output hashes. These prefixes are audit identifiers, not mathematical
certificates in themselves. The full archived programs contain the
rational endpoints; a compact representative verifier is included in
\cref{app:code} so this single-file edition contains a runnable example
without requiring external data.

\begin{longtable}{@{}P{0.24\textwidth}P{0.53\textwidth}P{0.17\textwidth}@{}}
\toprule
Verification family & Objects checked in the supplied implementation &
SHA-256 prefix\\
\midrule
\endhead
$\Ein$ area & Exact convolution, eight-equation system, denominator bounds,
and a factorial-series enclosure & \texttt{083e681137ea}\\
Gaussian mixtures & Five exact candidate classifications, two root
certificates, and reparametrized $E$-identities & \texttt{b5576f7986fa}\\
Student and beta & Four root intervals; 440 half-integral beta parameter
pairs; exact differential identities & \texttt{6057b3ca1d26}\\
Shortest beta intervals & 2,025 parameter triples; Bernstein signs and
coefficient identities; endpoint and length enclosures & \texttt{2e676a5f24d1}\\
Lemniscate length & Order-240 rational series, 964 coefficient equalities,
and a closed-form tail majorant & \texttt{a50ffd08d46e}\\
Normalized length & The same coefficient-square representation without
numerical division by an approximate $\pi$ & \texttt{2282da708d00}\\
Prime divisions & Three prime-leaf root enclosures and finite CM-type
checks, kept distinct from the general proof & \texttt{981cf8464013}\\
All cyclotomic divisions & 118 exponent projectors, 599 polynomial
remainders, 7,257 convolution checks, three root intervals &
\texttt{664439f0f33e}\\
Algebraic ratios & Two normalized integrals, a golden-ratio inverse, and
exact cyclotomic substitutions & \texttt{b0f6c3c92576}\\
Rational inputs & Seven numerical examples, good-prime choices, and
3,500 finite valuation checks & \texttt{58fda457e012}\\
Bounded degree & Three quadratic-input integrals; 18 exact quadratic
field calculations and 18,000 valuation inequalities & \texttt{50a3aea763a7}\\
First-kind Fermat beta & Three CDF values and three medians, 21 prime
pairs, and 243 polynomial-reduction identities & \texttt{1a8b36e134de}\\
Second-kind beta & Three CDF values and three medians, 27 prime pairs,
and first/second-kind character counts & \texttt{40eb9e56a1ab}\\
Gamma transfer & 58 polynomial tests, series identities, 484 Laurent
coefficient checks, and boundary valuation tests & \texttt{58179fa70170}\\
\bottomrule
\end{longtable}
The equal-output test establishes reproducibility of these implementations,
not independent mathematical proof: the same mistaken formula could
reproduce the same output. The core inequalities and algebraic identities
are therefore proved in the text rather than inferred from these counts.

\subsection{Representative certification scales}
The supplied verifiers certify the $\Ein$ area by an error radius less
than $1.76\times10^{-210}$ at order $160$. The two Gaussian medians use
intervals of width $10^{-66}$; the four classical quantile examples use
width $10^{-60}$. The beta HPD example and the principal arc-division
examples use intervals of width at most $10^{-70}$. At order $240$, the
length $L_3$ has a rational enclosure of width less than
$2.28\times10^{-83}$, and $L_3/\pi$ one of width less than
$8\times10^{-84}$. The rational-input, bounded-degree, and weighted-beta
CDF examples use rational enclosures of width below $10^{-100}$.
The two algebraic-position ratios from \cref{sec:ratios} are enclosed
within widths at most $2\times10^{-90}$; the golden-ratio inverse uses
width $10^{-65}$. Displayed decimals may use fewer digits and have wider
outward-rounded display intervals than the underlying rational
certificates. These two notions of width must not be interchanged.

\section{Synthesis and unresolved questions}\label{sec:conclusion}
The most reusable step in this manuscript is not an isolated special-value
identity. It is the identification of an arithmetic relation that can be
transported between valid settings. Differential translations distinguish
an antisymmetric iterated integral from products of its lower-level
primitives. Reparametrization puts Gaussian integrals inside the correct
$E$-function class. A closed square-root lift treats arc length without
assuming that sums or limits of transcendental numbers remain
transcendental. Cyclotomic period injectivity replaces simplicity where
a Jacobian decomposes. Torsion is the algebraic property that can be moved
from complex logarithms to local logarithms. Generalized Jacobians retain
finite polar information while projecting away additive directions only
when the period argument justifies it.

Several boundaries remain explicit. The three possible half-integral beta
pairs do not by themselves enumerate every integer parameter pair that
realizes the off-center exception. The algebraic-division exceptional set
for arbitrary composite exponent is finite, but is not shown here to be
empty for every algebraic input. The bounded-degree theorem has a large
explicit bound and does not classify all low-degree inputs at small
exponents. The polynomial-weight Fermat result does not automatically
extend to arbitrary rational weights, equal denominators, or every
composite denominator. Nor do the single-value results prove algebraic
independence among unrelated constants.

The lemniscate length and rational-division arguments address the precise
questions posed in \cite[Remarks 8.16 and 5.7]{Tamiozzo}, respectively,
in the form of proof manuscripts developed here. The fact that a question
appears in a 2019 source does not establish that it remained unanswered
until this edition. No assertion of first publication or community
acceptance is inferred from unsuccessful searches for a matching paper.
Independent review should focus on the geometric square-root lift,
the exact rational-character projection, the hypotheses of the torsion
classifications, and the passage from generalized logarithms to ordinary
abelian projections.

For Euler's constant the conclusion is different. The transfer identifies
coefficient locking and the arithmetic nature of the missing step, but
does not supply that step. The analytic identity
$\Psi_{r,0}(1)=e(\gamma-r)$ is exact; a theorem preventing its vanishing
for every algebraic $r$ is still needed. \Cref{conj:ED} and
\eqref{eq:missing-zero-lift} remain conditional targets, not inputs that
may be silently assumed. Consequently the manuscript ends with an audited
obstruction rather than a claimed proof of irrationality or transcendence
of $\gamma$.

\appendix
\section{Catalogue of explicit constants}\label{app:catalogue}
The entries below are exact numbers defined by the indicated integral or
inverse-function condition. Decimal strings are truncated for orientation.
Their arithmetic conclusions refer to the proofs in the cited sections,
not to the terminating decimal approximations. Letters used for local
examples have their definitions in those sections.

\begin{longtable}{@{}P{0.48\textwidth}P{0.36\textwidth}P{0.09\textwidth}@{}}
\toprule
Exact definition & Initial decimal digits & Section\\
\midrule
\endhead
$H_{12}(1)$, antisymmetric $\Ein$ area &
$0.066898077211666761862235\ldots$ & \ref{sec:ein}\\
Median of the asymmetric truncated Gaussian mixture on $[-1,1]$ &
$0.111102019216173792870452\ldots$ & \ref{sec:gaussian}\\
$J_1(m)=J_1(1)/2$, $0<m<1$ &
$0.392697625792947214547243\ldots$ & \ref{sec:gaussian}\\
$Q_3(39/40)$ for Student's $t$ law &
$3.182446305283709592723225\ldots$ & \ref{sec:classical-quantiles}\\
Median of $\operatorname{Beta}(3/2,5/2)$ &
$0.352452279877130813765599\ldots$ & \ref{sec:classical-quantiles}\\
$1/40$-quantile of $\operatorname{Beta}(3/2,5/2)$ &
$0.038747617785165160780701\ldots$ & \ref{sec:classical-quantiles}\\
$39/40$-quantile of the same law &
$0.823263902868742682562881\ldots$ & \ref{sec:classical-quantiles}\\
Left endpoint, $95\%$ shortest interval for $\operatorname{Beta}(3/2,5/2)$ &
$0.009556870873044942860219\ldots$ & \ref{sec:hpd}\\
Right endpoint of the same shortest interval &
$0.770757269257335428334836\ldots$ & \ref{sec:hpd}\\
Length of the same shortest interval &
$0.761200398384290485474616\ldots$ & \ref{sec:hpd}\\
Midpoint of the same shortest interval &
$0.390157070065190185597527\ldots$ & \ref{sec:hpd}\\
$L_3=\length\{|z^3-z|=1\}$ &
$6.464706327813477265623126\ldots$ & \ref{sec:lemniscate}\\
$L_3/\pi$ &
$2.057779935417939294967383\ldots$ & \ref{sec:lemniscate}\\
$s_{8,1/2}$, four-leaf half-arc midpoint &
$0.581372564921331418342397\ldots$ & \ref{sec:division}\\
$s_{10,1/2}$, five-leaf half-arc midpoint &
$0.566067311893637997084837\ldots$ & \ref{sec:division}\\
$s_{12,1/2}$, six-leaf half-arc midpoint &
$0.555521188888391003316451\ldots$ & \ref{sec:division}\\
$s_{18,1/2}$, nine-leaf half-arc midpoint &
$0.537498666292662410093407\ldots$ & \ref{sec:division}\\
$R_8(1/2)$ &
$0.429797068135111779727193\ldots$ & \ref{sec:ratios}\\
$R_{10}(1/2)$ &
$0.441595303662667131493307\ldots$ & \ref{sec:ratios}\\
$s_{8,(\sqrt5-1)/2}$ &
$0.716302719414143509845268\ldots$ & \ref{sec:ratios}\\
$R_{12}(1/2)$ &
$0.450017109099808036111667\ldots$ & \ref{sec:local}\\
$R_{18}(1/2)$ &
$0.465117308242266039709804\ldots$ & \ref{sec:local}\\
$R_{24}(1/2)$ &
$0.473209591762729788013422\ldots$ & \ref{sec:local}\\
$R_9(1/2)$ &
$0.436230919677225329059895\ldots$ & \ref{sec:local}\\
$R_{12}(2^{-1/4})=\mathrm I_{1/8}(1/12,1/2)$ &
$0.760658039134571678695270\ldots$ & \ref{sec:local}\\
$R_{12}(3-\sqrt7)$ &
$0.318832995555007621109640\ldots$ & \ref{sec:bounded}\\
$R_{18}(3-\sqrt7)$ &
$0.329534360277977105299820\ldots$ & \ref{sec:bounded}\\
$R_{24}(3-\sqrt7)$ &
$0.335267754547646648453111\ldots$ & \ref{sec:bounded}\\
$\mathrm I_{1/2}(1/3,1/5)$ &
$0.367139793949012364863668\ldots$ & \ref{sec:fermat}\\
Median of $\operatorname{Beta}(1/3,1/5)$ &
$0.804376357494993257372466\ldots$ & \ref{sec:fermat}\\
$\mathrm I_{1/2}(4/3,6/5)$ &
$0.460657366225158194044108\ldots$ & \ref{sec:fermat}\\
Median of $\operatorname{Beta}(4/3,6/5)$ &
$0.534154161604607302042768\ldots$ & \ref{sec:fermat}\\
CDF at $1/2$ for $a=1/3,b=1/5,W=1+t^2$ &
$0.248596392472207792663111\ldots$ & \ref{sec:fermat}\\
Median for $a=1/3,b=1/5,W=1+t^2$ &
$0.938730288889941676673362\ldots$ & \ref{sec:fermat}\\
$\mathrm I_{1/2}(2/3,2/5)$ &
$0.35389888866285687131\ldots$ & \ref{sec:fermat}\\
Median of $\operatorname{Beta}(2/3,2/5)$ &
$0.71685421441072752258\ldots$ & \ref{sec:fermat}\\
$\mathrm I_{1/2}(5/3,7/5)$ &
$0.43064354362164795497\ldots$ & \ref{sec:fermat}\\
Median of $\operatorname{Beta}(5/3,7/5)$ &
$0.553989284028517005618288\ldots$ & \ref{sec:fermat}\\
CDF at $1/2$ for $a=2/3,b=2/5,W=1+t^2$ &
$0.25205276478993036618\ldots$ & \ref{sec:fermat}\\
Median for $a=2/3,b=2/5,W=1+t^2$ &
$0.83890381650993691841\ldots$ & \ref{sec:fermat}\\
\bottomrule
\end{longtable}
No row in this catalogue is a claim about algebraic independence between
separate families. In particular the gamma constant itself is deliberately
not entered as a newly classified number.

\section{Concordance with the research notes}\label{app:concordance}
This single paper consolidates fourteen supplied notes. The table gives
English descriptive titles and their destination in the unified text.
These notes are internal research manuscripts, not external peer-reviewed
references; external inputs are listed separately in the bibliography.
Repetition has been removed by incorporating the strongest applicable
version while retaining earlier examples and explicit failures.

\begin{longtable}{@{}P{0.08\textwidth}P{0.66\textwidth}P{0.20\textwidth}@{}}
\toprule
Note & Subject of the supplied source & Unified location\\
\midrule
\endhead
N01 & Antisymmetric $\Ein$ iterated integrals and algebraic independence &
Section~\ref{sec:ein}\\
N02 & Truncated Gaussian mixtures and a finite algebraic-quantile test &
Section~\ref{sec:gaussian}\\
N03 & Student, half-integral beta, Jeffreys, and Clopper--Pearson quantiles &
Section~\ref{sec:classical-quantiles}\\
N04 & Arithmetic dichotomy for shortest beta probability intervals &
Section~\ref{sec:hpd}\\
N05 & Polynomial lemniscate lengths as closed periods &
Section~\ref{sec:lemniscate}\\
N06 & Prime-leaf rational divisions by simple-Jacobian rigidity &
Subsumed by Section~\ref{sec:division}\\
N07 & Normalized lemniscate length and logarithm separation &
Section~\ref{sec:lemniscate}\\
N08 & Cyclotomic proof for all integral leaf counts and exponents &
Section~\ref{sec:division}\\
N09 & Algebraic division ratios, cyclotomic conjugates, and torsion &
Section~\ref{sec:ratios}\\
N10 & Rational radii, Zsigmondy primes, and local nonvanishing &
Section~\ref{sec:local}\\
N11 & Height-independent bounds for algebraic inputs of bounded degree &
Section~\ref{sec:bounded}\\
N12 & First-kind Fermat beta integrals and polynomial-weight stability &
Section~\ref{sec:fermat}\\
N13 & Second-kind beta forms and generalized-Jacobian projection &
Section~\ref{sec:fermat}\\
N14 & Transfer audit for Euler's constant and the missing mixed-arithmetic
step & Section~\ref{sec:gamma}\\
\bottomrule
\end{longtable}
The N06 and N07 notes shared one computational archive, while the
associated total-length and normalized-length programs remain separate.
The gamma audit also recalled a September exploratory record on ten
approaches. That record is contextual rather than an additional source of
proved transcendence results. Its principal lessons---that finite
experiments do not prove impossibility, that denominator savings must
beat total arithmetic cost, and that nonsplitting alone does not determine
a comparison-matrix entry---are retained in Section~\ref{sec:gamma}.

\section{A compact exact-arithmetic verifier}\label{app:code}
The following self-contained Python listing uses only the standard library.
It verifies selected coefficient and valuation identities, then certifies
$R_{12}(1/2)$ in a rational interval of width below $10^{-100}$. The printed
60-place enclosure is intentionally wider than the internal interval.
It was executed in this form for this edition. It does not implement the
analytic subgroup theorem and is not a formal proof of transcendence.

\begin{lstlisting}
"""Exact rational checks; Python 3.10+, standard library only.
This code certifies a numerical interval, not transcendence.
"""
from fractions import Fraction as Q
from math import comb, factorial


def iroot(m: int, n: int) -> int:
    if m < 0 or n < 1:
        raise ValueError("Nonnegative radicand and positive degree required")
    lo, hi = 0, 1 << ((m.bit_length() + n - 1) // n)
    while lo + 1 < hi:
        mid = (lo + hi) // 2
        if mid ** n <= m:
            lo = mid
        else:
            hi = mid
    return hi if hi ** n <= m else lo


def root_bounds(q: Q, n: int, digits: int = 110):
    if q < 0:
        raise ValueError("Positive real root expected")
    scale = 10 ** digits
    m = iroot(q.numerator * scale ** n // q.denominator, n)
    return Q(m, scale), Q(m + 1, scale)


def beta_half(a: Q, b: Q, terms: int = 350):
    if not (0 < a < 1 and 0 < b < 1):
        raise ValueError("This positive-tail bound assumes 0 < a,b < 1")
    term, total = 1 / a, Q(0)
    for j in range(terms):
        total += term
        term *= (1 - b + j) * (a + j)
        term /= 2 * (j + 1) * (a + j + 1)
    lo, hi = root_bounds(Q(1, 2 ** a.numerator), a.denominator)
    # Consecutive terms have ratio < 1/2; term is first omitted.
    return lo * total, hi * (total + 2 * term)


def arc_bounds(n: int, s: Q, terms: int = 80):
    if n < 3 or not 0 < s < 1:
        raise ValueError("Need n >= 3 and 0 < s < 1")
    term, total, power = s, Q(0), s ** n
    for j in range(terms):
        total += term
        term *= Q(2 * j + 1, 2 * j + 2) * power
        term *= Q(n * j + 1, n * (j + 1) + 1)
    return total, total + term / (1 - power)


def ratio_bounds(n: int, s: Q):
    a, b = Q(1, n), Q(1, 2)
    p, q = beta_half(a, b), beta_half(b, a)
    den = ((p[0] + q[0]) / n, (p[1] + q[1]) / n)
    num = arc_bounds(n, s)
    assert den[0] > 0
    return num[0] / den[1], num[1] / den[0]


def decimal_floor(q: Q, digits: int = 60) -> str:
    scale = 10 ** digits
    m = q.numerator * scale // q.denominator
    return str(m // scale) + "." + str(m % scale).zfill(digits)


def check_coefficients():
    # Two independent expressions for coefficients of H_12(z).
    for n in range(2, 41):
        a = Q((-1) ** n, n) * sum(
            (Q(comb(n, k) * (2 ** k - 2 ** (n-k)), k)
             for k in range(1, n)), Q(0))
        ordinary = sum(
            (Q((-1) ** n * (2 ** k - 2 ** (n-k)),
               n * k * factorial(k) * factorial(n-k))
             for k in range(1, n)), Q(0))
        assert a == factorial(n) * ordinary
    # Universal local denominator bound on a finite test window.
    for ell in (5, 7, 11, 13):
        for j in range(1, 501):
            m, v = 2 * j + 1, 0
            while m % ell == 0:
                m //= ell
                v += 1
            assert (ell - 1) * v <= 2 * j


if __name__ == "__main__":
    check_coefficients()
    lo, hi = ratio_bounds(12, Q(1, 2))
    assert hi - lo < Q(1, 10 ** 100)
    scale = 10 ** 60
    out_lo = Q(lo.numerator * scale // lo.denominator, scale)
    out_hi = Q(-((-hi.numerator * scale) // hi.denominator), scale)
    assert out_lo <= lo < hi <= out_hi
    print("Lower:", decimal_floor(out_lo))
    print("Upper:", decimal_floor(out_hi))
    print("Exact interval width < 10^(-100).")
\end{lstlisting}
The output is
\begin{quote}\ttfamily\small
Lower: 0.450017109099808036111667642573989081152617667160799532175248\par
Upper: 0.450017109099808036111667642573989081152617667160799532175249\par
Exact interval width $<10^{-100}$.
\end{quote}
The internal width assertion and the wider decimal output concern different
intervals. The rational upper endpoint is rounded upward, and the lower
endpoint downward. The verification of finitely many local inequalities
is supplemental to the proof valid for every index in \cref{lem:local}.

\clearpage
\begin{thebibliography}{99}
\addcontentsline{toc}{section}{References}
\bibitem{Beukers}
F. Beukers, \emph{A refined version of the Siegel--Shidlovskii theorem},
Annals of Mathematics \textbf{163} (2006), 369--379.
\href{https://doi.org/10.4007/annals.2006.163.369}{doi:10.4007/annals.2006.163.369}.
Theorems 1.1 and 1.3 and Corollary 1.4.

\bibitem{Srinivasan}
V. Ravi Srinivasan, \emph{Extensions by Antiderivatives, Exponentials of
Integrals and by Iterated Logarithms}, preprint, 2008.
\url{https://arxiv.org/abs/0811.3004}. The Kolchin--Ostrowski theorem in
Section 2 is the form used here.

\bibitem{Baker}
A. Baker, \emph{Linear forms in the logarithms of algebraic numbers (III)},
Mathematika \textbf{14} (1967), 220--228.
\href{https://doi.org/10.1112/S0025579300003843}{doi:10.1112/S0025579300003843}.

\bibitem{HW}
A. Huber and G. W\"ustholz, \emph{Transcendence and Linear Relations of
1-Periods}, Cambridge Tracts in Mathematics 227, Cambridge University
Press, 2022. Author manuscript:
\url{https://arxiv.org/abs/1805.10104}.
Theorem 6.2, the generalized-Jacobian framework, and Corollary 13.13 are
used with their stated hypotheses.

\bibitem{Milne}
J. S. Milne, \emph{Algebraic Groups: The Theory of Group Schemes of Finite
Type over a Field}, Cambridge University Press, 2017. In particular the
Barsotti--Chevalley structure theorem.
\url{https://www.jmilne.org/math/Books/AGbook.pdf}.

\bibitem{FRLiouville}
S. Fischler and T. Rivoal, \emph{Values of $E$-functions are not Liouville
numbers}, Journal de l'\'Ecole polytechnique---Math\'ematiques
\textbf{11} (2024), 1--18.
\href{https://doi.org/10.5802/jep.249}{doi:10.5802/jep.249}.
\url{https://arxiv.org/abs/2301.01158}.
Theorem A restates the quantitative result used for the $\Ein$ example.

\bibitem{Tamiozzo}
M. Tamiozzo, \emph{Algebraicity of the division points of the trifolium
and related topics}, preprint, 2019.
\url{https://arxiv.org/abs/1912.00204}.
The arc-length formula in Section 3.4, the small-leaf results, and
Remarks 5.7 and 8.16 provide the historical problem statements discussed
here.

\bibitem{Orevkov}
S. Yu. Orevkov, \emph{Irreducibility of lemniscates}, preprint, 2018.
\url{https://arxiv.org/abs/1806.10377}.
The needed irreducibility argument is also given directly in the text.

\bibitem{CaroccaRodriguez}
A. Carocca and R. E. Rodr\'iguez, \emph{Jacobians with group actions and
rational idempotents}, Journal of Algebra \textbf{306} (2006), 322--343.
\url{https://arxiv.org/abs/math/0305328}.

\bibitem{BoxallGrant}
J. Boxall and D. Grant, \emph{Examples of torsion points on genus two
curves}, Transactions of the American Mathematical Society
\textbf{352} (2000), 4533--4555.
\href{https://doi.org/10.1090/S0002-9947-00-02368-0}{doi:10.1090/S0002-9947-00-02368-0}.
Theorem (0.3), the 22-point packet for $y^2=x^5+x$.

\bibitem{Arul}
V. Arul, \emph{Torsion points on Fermat quotients of the form
$y^n=x^d+1$}, preprint, 2019, revised version 2020.
\url{https://arxiv.org/abs/1910.14251}.
Definition 1.0.7, Proposition 2.4.4, Corollary 2.4.5, and Theorem 6.2.21.
The geometric torsion classification, not just rational-point torsion,
is the input required here.

\bibitem{Roessler}
D. Roessler, \emph{A note on the Manin--Mumford conjecture}, preprint,
2004. \url{https://arxiv.org/abs/math/0409083}.
Used for the torsion-coset formulation, not for an effective bound.

\bibitem{Zarhin}
Yu. G. Zarhin, \emph{$p$-adic abelian integrals and commutative Lie groups},
preprint, 1996. \url{https://arxiv.org/abs/alg-geom/9603006}.
Sections 1 and 3: functorial logarithms and local abelian integrals.

\bibitem{Coleman}
R. F. Coleman, \emph{Torsion points on curves and $p$-adic Abelian integrals},
Annals of Mathematics \textbf{121} (1985), 111--168.
\href{https://doi.org/10.2307/1971194}{doi:10.2307/1971194}.
Background for the local torsion strategy; the local series in this paper
is derived explicitly.

\bibitem{Avci}
\"O. Avc\i, \emph{Large Zsigmondy Primes}, preprint, 2020, revised
version 2023. \url{https://arxiv.org/abs/2011.06136}.
Only the classical Bang--Zsigmondy theorem recalled in the introduction
is needed here, not the paper's stronger large-prime result.

\bibitem{Zsigmondy}
K. Zsigmondy, \emph{Zur Theorie der Potenzreste}, Monatshefte f\"ur
Mathematik und Physik \textbf{3} (1892), 265--284.
\href{https://doi.org/10.1007/BF01692444}{doi:10.1007/BF01692444}.

\bibitem{Voutier}
P. M. Voutier, \emph{Primitive divisors of Lucas and Lehmer sequences, II},
Journal de Th\'eorie des Nombres de Bordeaux \textbf{8} (1996), 251--274.
\href{https://doi.org/10.5802/jtnb.168}{doi:10.5802/jtnb.168}.
Author version: \url{https://arxiv.org/abs/1211.3107}.
Theorem 2(i) is used; its ratio-degree parameter is distinguished from
the degree of an auxiliary ambient field.

\bibitem{Zannier}
U. Zannier, \emph{Hyperelliptic Continued Fractions and Generalized
Jacobians}, preprint, 2016. \url{https://arxiv.org/abs/1602.00934}.
Section 2.1.2 describes the modulus extension used in the second-kind
argument.

\bibitem{ACL}
D. Agostini, T. \"O. \c{C}elik, and J. B. Little,
\emph{Polynomial Analogs of Theta Functions and Rational Solutions of the
KP Equation}, International Mathematics Research Notices
\textbf{2024} (22), 13964--13987.
\href{https://doi.org/10.1093/imrn/rnae227}{doi:10.1093/imrn/rnae227}.
Sections 2.1 and 2.3 supply the Rosenlicht and Abel-map descriptions.

\bibitem{FischlerRivoalGevrey}
S. Fischler and T. Rivoal, \emph{Relations between values of arithmetic
Gevrey series, and applications to values of the Gamma function}, Journal
of Number Theory \textbf{261} (2024), 36--54.
\href{https://doi.org/10.1016/j.jnt.2024.02.016}{doi:10.1016/j.jnt.2024.02.016}.
\url{https://arxiv.org/abs/2301.13518}.
Conjectures 1 and 3, Definition 1, and Theorem 2 are explicitly separated
in the gamma audit.

\bibitem{RivoalSurvey}
T. Rivoal, \emph{Les $E$-fonctions et $G$-fonctions de Siegel},
Journ\'ees X-UPS 2019.
\href{https://doi.org/10.5802/xups.2019-03}{doi:10.5802/xups.2019-03}.
Theorem 5.11 and its discussion of $e$, $\Ein(1)$, and the
Euler--Gompertz alternative.

\bibitem{DhillonSaha}
S. Dhillon and E. Saha, \emph{On arithmetic nature of special values of
the incomplete beta function}, Acta Arithmetica \textbf{218} (2025),
273--284.
\href{https://doi.org/10.4064/aa240610-6-10}{doi:10.4064/aa240610-6-10}.
Related literature; a complete theorem-by-theorem priority comparison
with this paper has not been performed.

\bibitem{DLMF}
NIST Digital Library of Mathematical Functions, Chapters 5, 6, 7, and 8:
gamma, beta, exponential-integral, error-function, and incomplete-beta
identities. \url{https://dlmf.nist.gov/}.
Used for standard definitions and functional identities, not as a source
for the new applications assembled in this manuscript.

\bibitem{NIST}
NIST/SEMATECH, \emph{e-Handbook of Statistical Methods}, Sections
1.3.6.6.5, 1.3.6.6.17, and 1.3.6.7.2: $F$ distribution, beta distribution,
and Student's $t$ critical values.
\url{https://www.itl.nist.gov/div898/handbook/}.

\bibitem{Ratesci}
P. Laud, \emph{ratesci: Confidence Intervals for Comparisons of Binomial
or Poisson Rates}, CRAN package documentation, vignette ``1: Single
rate'' and reference page \texttt{jeffreysci}.
\url{https://cran.r-project.org/web/packages/ratesci/}.
The distinction between raw posterior quantiles and boundary-adjusted
interval endpoints is retained in the text.

\bibitem{Lourie}
N. Lourie, K. Cho, and H. He, \emph{Show Your Work with Confidence:
Confidence Bands for Tuning Curves}, preprint, 2023.
\url{https://arxiv.org/abs/2311.09480}.
Appendix C: the beta highest-density interval and equal-endpoint-density
condition.
\end{thebibliography}
\end{document}
